% =====================================================================
%  Física y modelado dinámico de un motor BLDC con acople magnético
%  y carga fluida — Guía progresiva desde las leyes fundamentales
%  Compilar con:  pdflatex fundamentos_fisicos_modelado.tex (x2)
%             o:  latexmk -pdf fundamentos_fisicos_modelado.tex
% =====================================================================
\documentclass[11pt,a4paper,oneside]{book}

\usepackage[utf8]{inputenc}
\usepackage[T1]{fontenc}
\usepackage[spanish,es-noshorthands,es-tabla]{babel}
\usepackage{lmodern}
\usepackage{microtype}
\usepackage[top=2.6cm,bottom=2.6cm,left=2.6cm,right=2.4cm,headheight=24pt]{geometry}
\usepackage{amsmath,amssymb}
\usepackage{cancel}
\usepackage{bm}
\usepackage{siunitx}
\sisetup{
  output-decimal-marker={,},
  group-separator={.},
  per-mode=symbol,
  detect-all
}
\usepackage{booktabs}
\usepackage{array}
\usepackage{enumitem}
\usepackage{xcolor}
\usepackage{fancyhdr}
\usepackage{tikz}
\usetikzlibrary{arrows.meta,positioning,calc,decorations.pathmorphing,patterns,shapes.geometric}
\usepackage{pgfplots}
\pgfplotsset{compat=1.16}
\usepackage[skins]{tcolorbox}
\usepackage[hidelinks,bookmarks=true,bookmarksopen=true]{hyperref}

% ---- Paleta sobria estilo libro de ingeniería ----
\definecolor{inkblue}{HTML}{1B3A5C}
\definecolor{tealaccent}{HTML}{196463}
\definecolor{amberwarn}{HTML}{9A5B13}
\definecolor{crimson}{HTML}{9B2226}
\definecolor{slatebg}{HTML}{F5F7FA}
\definecolor{tealbg}{HTML}{F1F8F7}
\definecolor{amberbg}{HTML}{FDF8F0}
\definecolor{quizbg}{HTML}{F3F5F9}
\definecolor{mutedgray}{HTML}{5C6773}

% ---- Encabezados y pies de página estilo libro ----
\pagestyle{fancy}
\fancyhf{}
\fancyhead[L]{\small\color{mutedgray}\nouppercase{\leftmark}}
\fancyhead[R]{\small\color{inkblue}\bfseries\thepage}
\renewcommand{\headrulewidth}{0.4pt}
\renewcommand{\headrule}{\hbox to\headwidth{\color{inkblue!25}\leaders\hrule height \headrulewidth\hfill}}

\fancypagestyle{plain}{
  \fancyhf{}
  \fancyfoot[C]{\small\color{inkblue}\bfseries\thepage}
  \renewcommand{\headrulewidth}{0pt}
}

% ---- Estilo limpio de títulos de capítulo (sin titlesec) ----
\makeatletter
\def\@makechapterhead#1{%
  \vspace*{10\p@}%
  {\parindent \z@ \raggedright \normalfont
    \ifnum \c@secnumdepth >\m@ne
      \if@mainmatter
        {\large\scshape\color{tealaccent}\@chapapp\space \thechapter}%
        \par\nobreak
        \vskip 6\p@
      \fi
    \fi
    \interlinepenalty\@M
    {\Huge \bfseries \color{inkblue} #1\par\nobreak}%
    \vskip 8\p@
    {\color{inkblue!35}\hrule height 1.2pt}%
    \vskip 24\p@
  }}
\def\@makeschapterhead#1{%
  \vspace*{10\p@}%
  {\parindent \z@ \raggedright \normalfont
    \interlinepenalty\@M
    {\Huge \bfseries \color{inkblue} #1\par\nobreak}%
    \vskip 8\p@
    {\color{inkblue!35}\hrule height 1.2pt}%
    \vskip 24\p@
  }}
\makeatother

% ---- Cajas pedagógicas ----
\newtcolorbox{ideafisica}[1][Idea física]{
  enhanced,
  colback=tealbg, colframe=tealaccent,
  boxrule=0pt, leftrule=3.5pt, arc=2pt,
  top=8pt, bottom=8pt, left=11pt, right=11pt,
  before skip=10pt, after skip=10pt,
  fonttitle=\bfseries\color{tealaccent},
  title={#1},
  attach title to upper={\par\smallskip}
}

\newtcolorbox{lectura}[1][Cómo se lee término a término]{
  enhanced,
  colback=slatebg, colframe=inkblue,
  boxrule=0pt, leftrule=3.5pt, arc=2pt,
  top=8pt, bottom=8pt, left=11pt, right=11pt,
  before skip=10pt, after skip=10pt,
  fonttitle=\bfseries\color{inkblue},
  title={#1},
  attach title to upper={\par\smallskip}
}

\newtcolorbox{unidades}[1][Verificación dimensional]{
  enhanced,
  colback=white, colframe=inkblue!40,
  boxrule=0.6pt, arc=2pt,
  top=6pt, bottom=6pt, left=10pt, right=10pt,
  before skip=8pt, after skip=10pt,
  fonttitle=\bfseries\small\color{inkblue},
  title={#1},
  attach title to upper={\quad}
}

\newtcolorbox{antidoto}[1][Qué dejamos afuera a propósito (antídoto contra la página en negro)]{
  enhanced,
  colback=amberbg, colframe=amberwarn,
  boxrule=0pt, leftrule=3.5pt, arc=2pt,
  top=8pt, bottom=8pt, left=11pt, right=11pt,
  before skip=12pt, after skip=10pt,
  fonttitle=\bfseries\color{amberwarn},
  title={#1},
  attach title to upper={\par\smallskip}
}

\newtcolorbox{lapizypapel}[1][Prueba de lápiz y papel (antes de seguir al próximo capítulo)]{
  enhanced,
  colback=quizbg, colframe=inkblue!75,
  boxrule=0.8pt, arc=3pt,
  top=8pt, bottom=8pt, left=11pt, right=11pt,
  before skip=12pt, after skip=14pt,
  fonttitle=\bfseries\color{white},
  colbacktitle=inkblue,
  title={#1}
}

\newtcolorbox{ecuacionclave}{
  enhanced,
  colback=slatebg, colframe=inkblue!45,
  boxrule=0.7pt, arc=2pt,
  top=5pt, bottom=5pt, left=8pt, right=8pt,
  before skip=8pt, after skip=8pt
}

% ---- Comandos matemáticos ----
\newcommand{\dd}{\mathrm{d}}
\newcommand{\sgn}{\operatorname{sgn}}
\newcommand{\sat}{\operatorname{sat}}
\setlength{\parskip}{0.45em}
\setlength{\parindent}{0pt}

% =====================================================================
\begin{document}
% =====================================================================

\frontmatter
\hypersetup{pageanchor=false}

% ---- Portada ----
\begin{titlepage}
  \centering
  \vspace*{1.6cm}
  {\small\scshape\color{tealaccent} Cuaderno de estudio --- Física y Control de Sistemas\par}
  \vspace{0.8cm}
  {\color{inkblue}\rule{\linewidth}{1.6pt}\par}
  \vspace{0.9cm}
  {\Huge\bfseries\color{inkblue}
    Física y modelado dinámico\\[0.35em]
    de un motor BLDC con acople magnético\\[0.35em]
    y carga fluida\par}
  \vspace{0.6cm}
  {\Large\color{mutedgray}\itshape
    Cómo construir el modelo desde cero con lápiz y papel:\\
    de las leyes elementales a la linealización y los límites de validez\par}
  \vspace{0.9cm}
  {\color{inkblue}\rule{\linewidth}{1.6pt}\par}
  \vspace{1.8cm}

  % Diagrama minimalista en portada
  \begin{tikzpicture}[
    node distance=1.22cm,
    box/.style={draw=inkblue!70, fill=slatebg, rounded corners=3pt,
                minimum width=2.15cm, minimum height=1.08cm, align=center,
                font=\scriptsize\color{inkblue}},
    >={Stealth[length=5pt]}
  ]
    \node[box] (e) {\textbf{Dominio}\\\textbf{eléctrico}\\$V,\,I$};
    \node[box, right=of e] (m1) {\textbf{Rotor}\\\textbf{del motor}\\$J_1,\,\omega_1$};
    \node[box, right=of m1, fill=amberbg, draw=amberwarn] (ac) {\textbf{Acople}\\\textbf{magnético}\\$T_c(\delta)$};
    \node[box, right=of ac] (m2) {\textbf{Agitador}\\\textbf{y fluido}\\$J_2,\,\omega_2$};
    \draw[->, thick, inkblue] (e) -- node[above, font=\scriptsize]{$T_e = K I$} (m1);
    \draw[->, thick, amberwarn] (m1) -- node[above, font=\scriptsize]{$\delta$} (ac);
    \draw[->, thick, amberwarn] (ac) -- node[above, font=\scriptsize]{$T_c$} (m2);
    \draw[->, dashed, tealaccent] (m1.south) -- ++(0,-0.65) -| node[pos=0.25, below, font=\scriptsize]{BEMF $e = K\omega_1$} (e.south);
  \end{tikzpicture}

  \vfill
  {\normalsize\color{mutedgray}
    Enfoque progresivo basado exclusivamente en leyes físicas comprensibles\\
    sin cajas negras ni parámetros sin sustento\par}
  \vspace{0.6cm}
  {\small\color{inkblue} Guía teórica de acompañamiento --- Control y Sistemas\par}
  \vspace*{0.6cm}
\end{titlepage}
\hypersetup{pageanchor=true}

% =====================================================================
\chapter*{Prefacio: cómo leer este libro}
\addcontentsline{toc}{chapter}{Prefacio: cómo leer este libro}
\markboth{Prefacio}{Prefacio}
% =====================================================================

Aprender a modelar un sistema electromecánico suele resultar intimidante cuando
uno abre un documento técnico y se encuentra, en la primera página, con una
matriz $4\times 4$, una función de transferencia de tercer orden y cinco
ecuaciones diferenciales ya armadas. Parece que hiciera falta memorizar fórmulas
de memoria. Sin embargo, ocurre exactamente lo contrario: \textbf{en ingeniería,
un buen modelo no se memoriza; se reconstruye en cinco minutos sobre una hoja en
blanco a partir de cuatro leyes físicas muy simples}.

Este libro fue escrito con un único propósito: acompañarte paso a paso en esa
construcción desde cero, olvidándonos por completo de detalles de programación,
registros de microcontroladores o bloques de simulador, para concentrarnos en
la \textbf{física real del sistema}: un motor sin escobillas (BLDC) que arrastra
una cápsula agitadora a través de un acople magnético dentro de un vaso con
fluido, regulado en velocidad sin sensor mecánico en el eje.

\section*{El «síndrome de la página en negro»}
Antes, al enfrentarse a un modelado, existía el miedo a la \emph{página en
blanco}: no saber por dónde empezar. Hoy existe el peligro opuesto, el
\textbf{síndrome de la página en negro}: llenar páginas con ecuaciones cada vez
más sofisticadas ---fricción dinámica de LuGre, tablas de turbulencia de tanques
industriales, modelos térmicos acoplados, cálculos de campo magnético en tres
dimensiones, efectos parásitos de conmutación--- simplemente porque es fácil
escribirlos o copiarlos.

Ese camino es una trampa por dos motivos:
\begin{enumerate}
  \item \textbf{Falsa sensación de seguridad:} si incluís un fenómeno cuyos
        parámetros no podés medir ni justificar, tu modelo empieza a simular un
        sistema imaginario que puede ser menos conservador que el real.
  \item \textbf{Pérdida de dominio conceptual:} si alguien te detiene con un
        lápiz y un papel y te pregunta \emph{«¿por qué este coeficiente lleva un
        signo menos y este otro un factor~2?»}, y no sabés deducirlo desde la
        física, el modelo dejó de ser tuyo.
\end{enumerate}

Por eso adoptamos aquí la regla de oro del modelado honesto: \textbf{preferimos
un modelo más simple, del que podamos explicar cada término, cada signo y cada
unidad con lápiz y papel, antes que un modelo recargado que no comprendamos por
completo}. Lo que no se modela se declara explícitamente, y siempre que hay
duda se toma una postura conservadora.

\section*{Las únicas cuatro leyes que vas a necesitar}
Todo el sistema ---desde la tensión de la fuente hasta la cápsula girando en el
líquido--- se deduce aplicando únicamente cuatro principios clásicos que ya
conocés de física general:
\begin{enumerate}
  \item \textbf{Ley de tensiones de Kirchhoff:} en una malla eléctrica, la
        tensión aplicada se reparte entre resistencia, inductancia y tensión
        inducida.
  \item \textbf{Ley de inducción de Faraday:} un imán moviéndose frente a una
        bobina induce una tensión proporcional a la velocidad de giro.
  \item \textbf{Conservación de la energía:} la potencia eléctrica que no se
        pierde en calor ni se guarda en el campo magnético \emph{tiene que}
        convertirse en potencia mecánica en el eje. Esta ley ata para siempre
        los dominios eléctrico y mecánico.
  \item \textbf{Segunda y tercera ley de Newton para rotación:} la suma de
        torques sobre cada cuerpo rígido iguala a su inercia por su aceleración
        angular ($\sum T = J\,\dot\omega$), y todo torque que un cuerpo ejerce
        sobre otro aparece en sentido opuesto sobre el primero.
\end{enumerate}

A esas cuatro leyes físicas sólo le sumaremos dos herramientas matemáticas de la
materia de control: la \textbf{aproximación lineal de Taylor} (reemplazar una
curva por su recta tangente cerca del punto de trabajo) y la
\textbf{transformada de Laplace} (para pasar de ecuaciones diferenciales a
funciones de transferencia y analizar un controlador PI).

\section*{Mapa del recorrido}
La figura siguiente resume la cadena física completa que vamos a ir armando
capítulo por capítulo. Cada caja corresponde a un capítulo del libro: primero
entendemos cada pieza aislada (Capítulos~\ref{cap:motor} a~\ref{cap:sensorless}),
luego las ensamblamos en un único modelo no lineal (Capítulo~\ref{cap:ensamble}),
lo linealizamos paso a paso (Capítulo~\ref{cap:linealizacion}) y finalmente
vemos cómo diseñamos un control lineal y en qué fronteras físicas exactas ese
modelo lineal deja de valer (Capítulo~\ref{cap:control_validez}).

\begin{figure}[htbp]
\centering
\begin{tikzpicture}[
  node distance=1.05cm and 0.95cm,
  blk/.style={draw=inkblue, fill=slatebg, rounded corners=3pt,
              align=center, minimum height=1.22cm, minimum width=2.25cm,
              font=\footnotesize, inner sep=4pt},
  nlblk/.style={draw=amberwarn, fill=amberbg, rounded corners=3pt,
                align=center, minimum height=1.22cm, minimum width=2.25cm,
                font=\footnotesize, inner sep=4pt},
  >={Stealth[length=5pt]}
]
  \node[blk] (pi) {\textbf{Controlador PI}\\+ saturación $V$\\(Cap.~\ref{cap:control_validez})};
  \node[blk, right=of pi] (mot) {\textbf{Motor eléctrico}\\$V \to I \to T_e$\\(Cap.~\ref{cap:motor})};
  \node[blk, right=of mot] (rot) {\textbf{Rotor ($J_1$)}\\Newton rotacional\\(Cap.~\ref{cap:mecanica})};
  \node[nlblk, right=of rot] (acp) {\textbf{Acople no lineal}\\$T_c(\delta)$\\(Cap.~\ref{cap:acople})};
  \node[blk, below=1.25cm of acp] (ag) {\textbf{Agitador ($J_2$)}\\+ fricción + fluido\\(Caps.~\ref{cap:mecanica} y~\ref{cap:cargas})};
  \node[nlblk, below=1.25cm of rot] (sen) {\textbf{Sensor BEMF}\\ceguera + retardo $1/n$\\(Cap.~\ref{cap:sensorless})};

  \node[left=0.45cm of pi, font=\footnotesize\bfseries, inkblue] (ref) {$\omega^\ast$};
  \draw[->, thick, inkblue] (ref) -- (pi);
  \draw[->, thick, inkblue] (pi) -- node[above, font=\scriptsize]{$d$} node[below, font=\tiny, mutedgray]{$[d_{\min},1]$} (mot);
  \draw[->, thick, inkblue] (mot) -- node[above, font=\scriptsize]{$T_e$} (rot);
  \draw[->, thick, amberwarn] (rot) -- node[above, font=\scriptsize]{$\theta_1,\omega_1$} (acp);
  \draw[->, thick, amberwarn] (acp) -- node[right, font=\scriptsize]{$T_c(\delta)$} (ag);
  \draw[->, dashed, inkblue] (rot.south) -- node[left, font=\scriptsize]{$\omega_1$} (sen.north);
  \draw[->, dashed, amberwarn] (sen.west) -| node[pos=0.28, below, font=\scriptsize]{$\hat\omega_1$ (atrasada o nula)} (pi.south);
  \draw[->, dotted, mutedgray] ($(acp.south)+(-0.45,0)$) -- ++(0,-0.45) -| node[pos=0.25, below, font=\tiny]{reacción $-T_c$} ($(rot.east)+(0,-0.25)$);
\end{tikzpicture}
\caption{Arquitectura física del sistema completo y guía de lectura por capítulos. En tono cálido se resaltan los bloques donde residen las no linealidades dominantes.}
\end{figure}

\tableofcontents

% =====================================================================
\mainmatter
% =====================================================================

% =====================================================================
\chapter{El motor eléctrico desde cero}\label{cap:motor}
% =====================================================================

\begin{ideafisica}
Un motor eléctrico no es una caja mágica: es simplemente una \textbf{bobina de
cobre colocada cerca de un imán}. Si le hacés pasar corriente, aparece un torque
que hace girar el imán. A la inversa, cuando el imán gira frente a la bobina,
induce una tensión que se opone a la que vos aplicás. Toda la dinámica eléctrica
del motor nace de esa doble relación.
\end{ideafisica}

Empecemos haciendo un experimento mental que podés comprobar con cualquier
motorcito de imán permanente sobre una mesa. Si lo desconectás de la fuente y
girás el eje con los dedos midiendo sus terminales con un voltímetro, aparece
una tensión. Cuanto más rápido lo girás, mayor es esa tensión. Ahora hacé lo
contrario: trabá el eje con la mano y conectale una tensión pequeña. Sentís un
esfuerzo de giro (un torque) proporcional a la corriente que circula por el
cable.

Esa simetría perfecta entre \emph{girar para generar tensión} y \emph{inyectar
corriente para generar torque} es el corazón de cualquier motor de imán
permanente, sea de corriente continua con escobillas o sin escobillas (BLDC).
Vamos a ponerla en ecuaciones despacio y sin saltarnos un solo paso.

% ---------------------------------------------------------------------
\section{Las dos leyes electromecánicas fundamentales}
% ---------------------------------------------------------------------

Consideremos primero un motor idealizado con un único circuito eléctrico
equivalente recorrido por una corriente $I(t)$, cuyo eje gira con velocidad
angular mecánica $\omega_1(t)$ (en \si{rad/s}).

\paragraph{1. Generación de tensión al girar (ley de Faraday).}
Cuando los imanes del rotor giran frente a las espiras del estator, el flujo
magnético que atraviesa las bobinas cambia con el tiempo. Según la ley de
Faraday, la tensión inducida es proporcional a la tasa de cambio del flujo, que
a su vez es proporcional a la velocidad angular $\omega_1$. A esa tensión
inducida la llamamos \textbf{fuerza contraelectromotriz} o \textbf{BEMF}
(\emph{Back Electromotive Force}), y la escribimos:
\begin{equation}
  e(t) = K\,\omega_1(t).
  \label{eq:bemf_gen}
\end{equation}
La constante $K$ depende sólo de la geometría del motor, la intensidad de sus
imanes y la cantidad de vueltas de cobre; por ahora no necesitamos saber qué
hay dentro del motor para usarla: nos dice \emph{cuántos volts genera el motor
por cada radián por segundo de giro}.

\paragraph{2. Generación de torque al circular corriente (fuerza de Lorentz).}
Cuando una corriente $I(t)$ circula por los conductores sumergidos en el campo
magnético fijo de los imanes permanentes, cada conductor siente una fuerza de
Lorentz proporcional al campo magnético y a la corriente. Como el campo de los
imanes es constante y el radio del rotor es fijo, el torque electromagnético
resultante es directamente proporcional a la corriente:
\begin{equation}
  T_e(t) = K_t\,I(t).
  \label{eq:torque_gen}
\end{equation}

En principio escribimos dos símbolos distintos, $K$ en~\eqref{eq:bemf_gen} y
$K_t$ en~\eqref{eq:torque_gen}, porque vienen de dos fenómenos distintos
(Faraday y Lorentz). Pero en física dos dominios acoplados nunca son
independientes: están unidos por la \textbf{conservación de la energía}.

% ---------------------------------------------------------------------
\section{Conservación de la energía: por qué \texorpdfstring{$K_t$ y $K$}{Kt y K} son la misma constante}
% ---------------------------------------------------------------------

Imaginemos nuestro motor ideal (sin pérdidas internas en la conversión entre
electricidad y movimiento) girando a velocidad $\omega_1$ mientras circula una
corriente $I$:
\begin{itemize}
  \item Desde el lado \textbf{eléctrico}, la fuente debe vencer la tensión
        inducida $e$ empujando la corriente $I$. La potencia eléctrica que se
        transfiere al mecanismo de conversión es:
        \[
          P_{\text{em}} = e\cdot I = \big(K\,\omega_1\big)\,I.
        \]
  \item Desde el lado \textbf{mecánico}, el eje recibe un torque $T_e$ mientras
        gira a velocidad angular $\omega_1$. La potencia mecánica entregada al
        rotor es:
        \[
          P_{\text{mec}} = T_e\cdot \omega_1 = \big(K_t\,I\big)\,\omega_1.
        \]
\end{itemize}
Si no se crea ni se destruye energía en el entrehierro magnético, ambas
potencias deben ser idénticas para \emph{cualquier} corriente $I$ y
\emph{cualquier} velocidad $\omega_1$:
\begin{equation}
  P_{\text{em}} = P_{\text{mec}}
  \quad\Longrightarrow\quad
  K\,\omega_1\,I = K_t\,I\,\omega_1
  \quad\Longrightarrow\quad
  \boxed{\;K_t = K\;.}
  \label{eq:kt_igual_k}
\end{equation}

\begin{unidades}
Comprobemos que las unidades de $K$ y $K_t$ en el Sistema Internacional son
exactamente las mismas. Recordemos que $\SI{1}{V} = \SI{1}{J/C}$,
$\SI{1}{J} = \SI{1}{N.m}$, $\SI{1}{A} = \SI{1}{C/s}$ y que el radián es
adimensional:
\[
  [K] = \frac{\si{V}}{\si{rad/s}}
      = \frac{\si{J/C}}{\si{s^{-1}}}
      = \frac{\si{N.m}}{\si{C/s}}
      = \frac{\si{N.m}}{\si{A}}
      = [K_t].
\]
\textbf{Cuidado práctico:} la igualdad numérica $K_t = K$ vale \emph{sólo} en
unidades SI (\si{rad/s}, \si{V}, \si{N.m}, \si{A}). Si expresás la velocidad en
rpm o el voltaje en \si{mV/rpm}, aparece el factor de conversión $\pi/30$, y
olvidar esa conversión al pasar de un dominio a otro es uno de los errores más
graves de modelado. Por eso en física siempre derivamos en \si{rad/s} y recién
al final convertimos a rpm usando $n = \frac{30}{\pi}\omega_1$.
\end{unidades}

% ---------------------------------------------------------------------
\section{La ecuación eléctrica del motor paso a paso}
% ---------------------------------------------------------------------

Un bobinado real de cobre no sólo tiene tensión inducida $e(t)$: también opone
resistencia eléctrica al paso de corriente (calentándose por efecto Joule) y
almacena energía en su propio campo magnético (autoinductancia). Podemos
representar entonces el circuito eléctrico del motor como tres elementos en
serie conectados a la tensión aplicada $V(t)$, tal como se muestra en la
Figura~\ref{fig:circuito_motor}.

\begin{figure}[htbp]
\centering
\begin{tikzpicture}[
  >={Stealth[length=5pt]},
  thick
]
  % Fuente V(t)
  \draw (0,0) -- (0,1.0);
  \draw (0,1.5) circle (0.5cm);
  \node at (0,1.68) {\scriptsize $+$};
  \node at (0,1.30) {\scriptsize $-$};
  \node[left=0.6cm] at (0,1.5) {$V(t)$};
  \draw (0,2.0) -- (0,3.0) -- (1.3,3.0);

  % Flecha de corriente
  \draw[->, crimson, very thick] (0.45,3.25) -- (1.25,3.25) node[midway, above, font=\small]{$I(t)$};

  % Resistencia R_eq
  \draw[inkblue] (1.3,3.0) -- (1.6,3.0)
    -- (1.72,3.22) -- (1.96,2.78) -- (2.20,3.22) -- (2.44,2.78)
    -- (2.68,3.22) -- (2.92,2.78) -- (3.04,3.0) -- (3.5,3.0);
  \node[above=0.3cm, inkblue] at (2.35,3.0) {$R_{\text{eq}}$};

  % Inductancia L_eq
  \draw[inkblue] (3.5,3.0) -- (3.8,3.0)
    arc (180:0:0.18) arc (180:0:0.18) arc (180:0:0.18) arc (180:0:0.18)
    -- (5.8,3.0) -- (6.8,3.0) -- (6.8,2.15);
  \node[above=0.3cm, inkblue] at (4.55,3.0) {$L_{\text{eq}}$};

  % BEMF e(t) = K omega_1
  \draw[tealaccent, fill=tealbg] (6.8,1.5) circle (0.65cm);
  \node[tealaccent, font=\small\bfseries] at (6.8,1.5) {M};
  \node[left=0.72cm, tealaccent, font=\scriptsize] at (6.8,1.9) {$+$};
  \node[left=0.72cm, tealaccent, font=\scriptsize] at (6.8,1.1) {$-$};
  \node[left=1.12cm, align=right, tealaccent, font=\small] at (6.8,1.5) {$e(t) = K\,\omega_1$\\(BEMF)};

  % Retorno
  \draw (6.8,0.85) -- (6.8,0) -- (0,0);

  % Eje mecánico saliendo a la derecha
  \draw[double distance=2pt, inkblue] (7.45,1.5) -- (9.35,1.5);
  \draw[->, inkblue, thick] (8.55,1.95) arc (60:-60:0.45cm);
  \node[above, inkblue, font=\small] at (8.85,2.02) {$\omega_1,\, T_e$};
\end{tikzpicture}
\caption{Circuito eléctrico equivalente de un motor de imán permanente. La tensión aplicada $V(t)$ se equilibra con la caída óhmica, la autoinductancia y la BEMF generada por el giro.}
\label{fig:circuito_motor}
\end{figure}

Aplicando la \textbf{ley de tensiones de Kirchhoff} a la malla de la
Figura~\ref{fig:circuito_motor} (la suma de caídas iguala a la tensión aplicada):
\begin{ecuacionclave}
\begin{equation}
  V(t) = R_{\text{eq}}\,I(t) + L_{\text{eq}}\,\frac{\dd I(t)}{\dd t} + K\,\omega_1(t),
  \qquad
  T_e(t) = K\,I(t).
  \label{eq:motor_eq_general}
\end{equation}
\end{ecuacionclave}

\begin{lectura}
Imaginemos la tensión $V(t)$ como el «presupuesto de empuje eléctrico» que
entregamos al motor. Ese presupuesto se gasta en tres cosas muy concretas:
\begin{itemize}
  \item $R_{\text{eq}}\,I$: lo que se pierde en vencer la resistencia del cobre
        (se transforma en calor, $P_{\text{Joule}} = R_{\text{eq}} I^2$).
  \item $L_{\text{eq}}\,\dot I$: lo que cuesta \emph{cambiar} la corriente
        rápidamente (la inductancia actúa como una «inercia eléctrica» que
        almacena energía magnética $E_{\text{mag}} = \tfrac{1}{2}L_{\text{eq}}I^2$).
  \item $K\,\omega_1$: la oposición generada por el propio movimiento del rotor
        (BEMF). Cuanto más rápido gira el motor, más «empuja hacia atrás» su
        tensión inducida, reduciendo la corriente automáticamente.
\end{itemize}
Podemos verificar el balance de potencia multiplicando la primera ecuación
de~\eqref{eq:motor_eq_general} por $I$:
\[
  \underbrace{V\,I}_{\text{potencia fuente}}
  = \underbrace{R_{\text{eq}}\,I^2}_{\text{calor}}
  + \underbrace{\frac{\dd}{\dd t}\!\left(\frac{1}{2}L_{\text{eq}} I^2\right)}_{\text{energía magnética}}
  + \underbrace{T_e\,\omega_1}_{\text{potencia al eje}}.
\]
Cada término tiene un destino físico exacto; nada sobra ni falta.
\end{lectura}

% ---------------------------------------------------------------------
\section{¿Por qué un motor BLDC trifásico obedece esta misma ecuación?}
% ---------------------------------------------------------------------

Hasta aquí hablamos de un motor de corriente continua con dos bornes. Pero
nuestro agitador usa un \textbf{motor BLDC trifásico} (sin escobillas, como el
de un disco rígido). ¿Hace falta entonces modelar tres fases, seis transistores
y ángulos de conmutación para entender su dinámica de velocidad? La respuesta es
\textbf{no}, y entender \emph{por qué} es clave para no complicarse de más.

\subsection{Ángulo mecánico y ángulo eléctrico}
En un motor BLDC los imanes permanentes están pegados al rotor y las tres
bobinas ($a, b, c$, conectadas en estrella con un punto neutro interno $n$)
están fijas en el estator. Un motor puede tener $p = 1$ par de polos magnéticos
(un Norte y un Sur en el rotor), o bien $p = 2, 3, \dots$ pares de polos
distribuidos alrededor de la circunferencia.

Cada vez que el rotor gira una vuelta mecánica completa ($\theta_1 = 2\pi$),
frente a una misma bobina pasan $p$ polos Norte y $p$ polos Sur. Es decir, el
patrón electromagnético se repite $p$ veces por vuelta mecánica. Por eso
definimos el \textbf{ángulo eléctrico} y la \textbf{velocidad eléctrica}:
\begin{equation}
  \theta_e = p\,\theta_1, \qquad \omega_e = p\,\omega_1, \qquad f_e = \frac{\omega_e}{2\pi} = \frac{p\,\omega_1}{2\pi} = \frac{p\,n}{60}.
  \label{eq:angulo_electrico}
\end{equation}
Por ejemplo, si $p = 2$ pares de polos y el rotor da una vuelta mecánica por
segundo ($n = \SI{60}{rpm}$), cada bobina ve dos ciclos magnéticos completos por
segundo ($f_e = \SI{2}{Hz}$).

\subsection{Conmutación en seis pasos (\emph{six-step}): dos fases en serie y una libre}
En un motor DC tradicional, las escobillas mecánicas se encargan de ir
cambiando qué bobinas reciben corriente a medida que el rotor gira, para que el
campo del estator siempre esté a \SI{90}{\degree} eléctricos del rotor (torque
máximo). En un BLDC operado en \textbf{conmutación trapezoidal de seis pasos},
esa misma tarea la hacen seis interruptores electrónicos (un inversor
trifásico):
\begin{quote}
  \emph{En cada sector de \SI{60}{\degree} eléctricos, el inversor conecta una
  fase al positivo de la fuente, otra fase al negativo y deja la tercera fase
  completamente desconectada (flotante).}
\end{quote}

Miremos la Figura~\ref{fig:bldc_estrella} en un sector cualquiera, por ejemplo
cuando conducen la fase~$a$ (hacia el neutro) y la fase~$b$ (desde el neutro),
mientras la fase~$c$ queda abierta ($i_c = 0$).

\begin{figure}[htbp]
\centering
\begin{tikzpicture}[
  >={Stealth[length=5pt]},
  thick
]
  % Neutro
  \coordinate (N) at (0,0);
  \fill[inkblue] (N) circle (2.2pt) node[below left=2pt, font=\small]{$n$};

  % Fase A (izquierda-arriba)
  \draw[inkblue, very thick] (-3.2,1.6) -- (-1.6,0.8) -- (N);
  \node[fill=slatebg, draw=inkblue, rounded corners=2pt, font=\scriptsize, inner sep=3pt] at (-2.0,1.0) {$R,\, L_s,\, e_a$};
  \node[left, inkblue, font=\small\bfseries] at (-3.2,1.6) {$+V$ (Fase $a$)};
  \draw[->, crimson, very thick] (-2.9,1.8) -- (-2.1,1.4) node[midway, above right, font=\scriptsize]{$I$};

  % Fase B (derecha-arriba)
  \draw[inkblue, very thick] (N) -- (1.6,0.8) -- (3.2,1.6);
  \node[fill=slatebg, draw=inkblue, rounded corners=2pt, font=\scriptsize, inner sep=3pt] at (2.0,1.0) {$R,\, L_s,\, e_b$};
  \node[right, inkblue, font=\small\bfseries] at (3.2,1.6) {$0\text{ V}$ (Fase $b$)};
  \draw[->, crimson, very thick] (2.1,1.4) -- (2.9,1.8) node[midway, above left, font=\scriptsize]{$I$};

  % Fase C (abajo, flotante)
  \draw[mutedgray, dashed, thick] (N) -- (0,-1.85);
  \node[fill=white, draw=mutedgray, dashed, rounded corners=2pt, font=\scriptsize, inner sep=3pt, text=mutedgray] at (0,-1.0) {$R,\, L_s,\, e_c$};
  \node[below, tealaccent, font=\small] at (0,-1.85) {Fase $c$ flotante ($i_c = 0$, sirve para medir BEMF)};
\end{tikzpicture}
\caption{En conmutación de seis pasos siempre hay dos fases en serie llevando la misma corriente $I$ y una tercera fase abierta ($i_c=0$).}
\label{fig:bldc_estrella}
\end{figure}

Hagamos las cuentas directamente sobre la Figura~\ref{fig:bldc_estrella}. Como
$i_c = 0$, por conservación de carga en el nodo neutro $n$ la corriente que
entra por la fase~$a$ debe salir íntegra por la fase~$b$:
\[
  i_a = +I, \qquad i_b = -I, \qquad i_c = 0.
\]
Ahora sumemos las tres propiedades de las dos fases en serie:
\begin{enumerate}
  \item \textbf{Resistencia equivalente:} la corriente atraviesa dos bobinas de
        resistencia por fase $R$ en serie $\implies$
        \[
          R_{\text{eq}} = R + R = 2R.
        \]
  \item \textbf{Inductancia equivalente:} en cada fase, el flujo magnético
        propio vale $L\,i_k$ y además el acoplamiento mutuo con la otra fase
        activa aporta $M\,i_{\text{otra}}$. Como $i_b = -I$, el flujo en la
        fase~$a$ es $L I + M(-I) = (L - M)I \equiv L_s I$ (a $L_s = L - M$ se la
        llama \emph{inductancia síncrona por fase}). Sumando las dos fases en
        serie resulta:
        \[
          L_{\text{eq}} = 2 L_s.
        \]
        Lo más práctico es que, si apoyás las puntas de un multímetro o medidor
        LCR entre dos terminales cualesquiera del motor, medís \emph{directamente}
        $R_{LL} = 2R$ y $L_{LL} = 2L_s$.
  \item \textbf{BEMF y constante de torque:} si la conmutación está sincronizada
        con el rotor, mientras conducen $a$ y $b$ sus tensiones inducidas están
        en la meseta de signo opuesto: $e_a = +k_e\,\omega_1$ y
        $e_b = -k_e\,\omega_1$, donde $k_e$ es la constante de BEMF \emph{por
        fase}. La tensión inducida total entre terminales es:
        \[
          e_{ab} = e_a - e_b = (+k_e\omega_1) - (-k_e\omega_1) = 2k_e\,\omega_1.
        \]
        Por conservación de energía ($P_{\text{em}} = e_{ab}\,I = 2k_e\omega_1 I = T_e\omega_1$),
        la constante de torque equivalente queda fijada sin ambigüedad:
        \[
          K = k_t = 2k_e \qquad\Longrightarrow\qquad T_e = 2k_e\,I.
        \]
\end{enumerate}

Reemplazando $R_{\text{eq}}=2R$, $L_{\text{eq}}=2L_s$ y $K=2k_e$ en nuestra
ecuación general~\eqref{eq:motor_eq_general}, obtenemos exactamente el modelo
eléctrico del BLDC en conmutación de seis pasos:
\begin{ecuacionclave}
\begin{equation}
  V = 2R\,I + 2L_s\,\frac{\dd I}{\dd t} + 2k_e\,\omega_1,
  \qquad
  T_e = 2k_e\,I.
  \label{eq:bldc_sixstep}
\end{equation}
\end{ecuacionclave}

Observá que \emph{no hubo que adivinar ningún factor 2}: cada 2 proviene
físicamente de que siempre conducen exactamente dos fases en serie. Y si
preferís escribir $R_{LL}$, $L_{LL}$ y $K$ (medidos de línea a línea), la
ecuación es indistinguible de un motor DC clásico.

% ---------------------------------------------------------------------
\section{Modulación PWM y la primera no linealidad: saturación de tensión}
% ---------------------------------------------------------------------

En la práctica no variamos la tensión $V$ con una fuente analógica regulable,
sino con una fuente fija de tensión de bus $V_{dc}$ (por ejemplo \SI{12}{V})
conmutada rápidamente mediante \textbf{modulación por ancho de pulso (PWM)} a
una frecuencia $f_{\text{PWM}}$ del orden de decenas de kilohercios (típicamente
15--\SI{20}{kHz}).

Durante una fracción $d \in [0, 1]$ de cada período PWM (el \textbf{ciclo de
trabajo} o \emph{duty cycle}), el puente aplica $V_{dc}$; durante el resto
$(1-d)$, la corriente circula por el camino de rueda libre con tensión aplicada
nula. Como el período del PWM ($T_{\text{PWM}} \sim \SI{50}{\micro s}$) es mucho
más rápido que las constantes de tiempo eléctrica y mecánica del motor, las
bobinas actúan como un filtro pasabajos natural que «promedia» los pulsos y ve
una tensión efectiva continua:
\begin{equation}
  V(t) = d(t)\,V_{dc}.
  \label{eq:pwm_avg}
\end{equation}

Sin embargo, el sentido común nos impone inmediatamente una restricción física
ineludible: \textbf{el inversor jamás puede aplicar más tensión que la que tiene
la fuente}. En operación unidireccional a lazo cerrado, además suele imponerse un
piso técnico $d_{\min} > 0$ (p.\,ej.\ $d_{\min} = 0{,}05$) para mantener
actividad mínima de conmutación. Así aparece nuestra \textbf{primera no
linealidad dominante}:
\begin{equation}
  \boxed{\;d \in [d_{\min},\, 1] \quad\Longleftrightarrow\quad V \in [d_{\min}V_{dc},\, V_{dc}]\;.}
  \label{eq:sat_tension}
\end{equation}
Un controlador lineal ideal podría pedir $d = 2{,}5$ (o sea \SI{30}{V} a una
fuente de \SI{12}{V}) frente a un escalón grande de consigna o una traba. La
planta real simplemente se queda pegada al techo $V_{dc}$. Omitir esta
saturación sería un error de sentido común grave: haría creer que el motor
dispone de torque y aceleración infinitos.

% ---------------------------------------------------------------------
\section{Intuición de régimen permanente: curva torque--velocidad y amortiguamiento eléctrico}
% ---------------------------------------------------------------------

Antes de pasar a la mecánica, detengámonos un momento en~\eqref{eq:bldc_sixstep}
cuando la corriente ya se estabilizó ($\dd I/\dd t = 0$). Despejando $I$ y
sustituyendo en $T_e = 2k_e I$:
\begin{align}
  I &= \frac{d\,V_{dc} - 2k_e\,\omega_1}{2R}, \\[4pt]
  T_e(\omega_1) &= 2k_e\,I
                 = \underbrace{\frac{k_e}{R}\,d\,V_{dc}}_{T_{\text{arranque}}(d)}
                 - \underbrace{\frac{2k_e^2}{R}}_{B_e}\,\omega_1.
  \label{eq:curva_tw}
\end{align}

Esta recta, dibujada en la Figura~\ref{fig:curva_tw}, encierra casi toda la
intuición física de cómo «siente» un motor eléctrico la carga:
\begin{itemize}
  \item \textbf{A rotor trabado ($\omega_1 = 0$):} no hay BEMF ($e = 0$), la
        corriente sólo está limitada por la resistencia del cobre
        ($I_{\text{traba}} = d\,V_{dc}/(2R)$) y el motor entrega su \textbf{torque
        máximo} para ese ciclo de trabajo, $T_{\text{traba}} = \frac{k_e}{R}d\,V_{dc}$.
  \item \textbf{En vacío ideal ($T_e = 0$, $I = 0$):} el motor acelera hasta que
        su propia BEMF iguala exactamente a la tensión aplicada ($2k_e\omega_1 = d\,V_{dc}$),
        alcanzando la \textbf{velocidad de vacío}:
        \[
          \omega_{\text{vacío}}(d) = \frac{d\,V_{dc}}{2k_e}.
        \]
  \item \textbf{Amortiguamiento eléctrico $B_e = 2k_e^2/R$:} mirá el segundo
        término de~\eqref{eq:curva_tw}, $-B_e\,\omega_1$. Tiene exactamente la
        forma matemática de una fricción viscosa proporcional a la velocidad.
        Físicamente, aunque no haya rozamiento mecánico, la BEMF hace que si la
        velocidad sube, el torque baje automáticamente; y si la velocidad cae
        por una carga, la BEMF cae, entra más corriente y el torque sube solo.
        El motor eléctrico se \textbf{auto-estabiliza} en velocidad!
\end{itemize}

\begin{figure}[htbp]
\centering
\begin{tikzpicture}
\begin{axis}[
  width=10.8cm, height=5.8cm,
  axis lines=left,
  xlabel={Velocidad angular $\omega_1$},
  ylabel={Torque electromagnético $T_e$},
  xmin=0, xmax=1.15, ymin=0, ymax=1.15,
  xtick={0.5, 1.0},
  xticklabels={$\frac{0{,}5\,V_{dc}}{2k_e}$, $\frac{V_{dc}}{2k_e}$},
  ytick={0.5, 1.0},
  yticklabels={$0{,}5\,\frac{k_e V_{dc}}{R}$, $\frac{k_e V_{dc}}{R}$},
  tick label style={font=\footnotesize},
  label style={font=\small},
  clip=false
]
  % Zona inalcanzable por saturación d > 1
  \fill[crimson!8] (axis cs:0,1.0) -- (axis cs:1.0,0) -- (axis cs:1.15,0) -- (axis cs:1.15,1.15) -- (axis cs:0,1.15) -- cycle;
  \node[crimson, font=\scriptsize, align=center] at (axis cs:0.86,0.84) {Zona prohibida\\por saturación ($d>1$)};

  % Rectas para d = 1 y d = 0.5
  \addplot[inkblue, very thick, domain=0:1] {1 - x};
  \node[inkblue, font=\footnotesize, rotate=-24] at (axis cs:0.25,0.83) {$d = 1$ (techo $V_{dc}$)};

  \addplot[tealaccent, thick, dashed, domain=0:0.5] {0.5 - x};
  \node[tealaccent, font=\footnotesize, rotate=-24] at (axis cs:0.22,0.34) {$d = 0{,}5$};

  % Anotación pendiente
  \draw[inkblue!75, thin] (axis cs:0.46,0.54) -- (axis cs:0.66,0.54) -- (axis cs:0.66,0.34);
  \node[right, inkblue, font=\scriptsize] at (axis cs:0.68,0.45) {pendiente $-B_e = -\dfrac{2k_e^2}{R}$};
\end{axis}
\end{tikzpicture}
\caption{Característica estática torque--velocidad del motor eléctrico para dos ciclos de trabajo. La recta $d=1$ marca el límite físico duro impuesto por la tensión de bus $V_{dc}$.}
\label{fig:curva_tw}
\end{figure}

% ---------------------------------------------------------------------
\section{Escalas de tiempo: cuándo la inductancia es «rápida»}
% ---------------------------------------------------------------------

Si aplicamos un escalón de tensión con el rotor quieto ($\omega_1=0$), la
ecuación $2L_s\dot I + 2R I = V$ es una ecuación diferencial de primer orden
cuya corriente alcanza el régimen en una \textbf{constante de tiempo eléctrica}:
\begin{equation}
  \tau_e = \frac{L_{\text{eq}}}{R_{\text{eq}}} = \frac{2L_s}{2R} = \frac{L_s}{R}.
  \label{eq:tau_e}
\end{equation}
En motores pequeños como el de un disco rígido, $R$ suele estar en el orden de
unos pocos ohms y $L_s$ en el orden de fracciones de milihenry o pocos
milihenrys, por lo que $\tau_e$ está típicamente entre \SI{0,2}{ms} y
\SI{1}{ms}. En cambio, como veremos en los próximos capítulos, acelerar la masa
del rotor y del agitador toma decenas o cientos de milisegundos
($\tau_{\text{mec}} \sim 50\text{--}\SI{200}{ms}$).

Esa separación de escalas ($\tau_e \ll \tau_{\text{mec}}$) es una bendición:
\begin{itemize}
  \item en el modelo completo de simulación podemos conservar $L_s$ sin
        esfuerzo (es sólo un integrador más);
  \item pero cuando hagamos \textbf{análisis a mano con lápiz y papel}
        (Capítulo~\ref{cap:linealizacion}), podremos suponer que la corriente se
        acomoda casi instantáneamente ($2L_s\dot I \approx 0$) y usar
        directamente~\eqref{eq:curva_tw}, simplificando enormemente las cuentas
        sin perder precisión en la dinámica lenta de velocidad.
\end{itemize}

\begin{antidoto}
En este capítulo dejamos afuera deliberadamente:
\begin{itemize}
  \item \textbf{Conmutación transistor por transistor y rizado PWM:} ocurre a
        \SI{18}{kHz}, órdenes de magnitud más rápido que el movimiento mecánico.
  \item \textbf{Torque de \emph{cogging} (efecto de ranuras del estator):} es un
        rizado magnético periódico sin corriente que requeriría medir armónicos
        espaciales con torquímetro; a velocidades de giro continuas su promedio
        en una vuelta es cero.
  \item \textbf{Dinámica térmica del bobinado ($R(T)$):} requeriría identificar
        resistencias y capacitancias térmicas que no medimos; operar lejos de la
        corriente límite nos mantiene conservadores sin inventar constantes.
  \item \textbf{Forma exacta no ideal de la BEMF:} aunque un motor real tenga
        una BEMF intermedia entre trapezoidal y senoidal, en operación six-step
        el modelo de constante media efectiva $2k_e$ captura fielmente el balance
        energético global.
\end{itemize}
\end{antidoto}

\begin{lapizypapel}
\begin{enumerate}[leftmargin=1.5em]
  \item Dibujá el circuito equivalente del motor eléctrico y escribí su ecuación
        diferencial explicando el significado físico y las unidades de cada uno
        de los cuatro términos.
  \item Demostrá en tres líneas, usando conservación de energía, por qué en
        unidades SI la constante de torque y la de BEMF deben ser iguales
        ($K_t = K$). ¿Qué pasa si alguien mezcla rpm con \si{N.m/A}?
  \item Explicá por qué en un BLDC trifásico con conmutación six-step aparecen
        los factores~2 en $2R$, $2L_s$ y $2k_e$.
  \item ¿Qué significa físicamente el término $- \frac{2k_e^2}{R}\omega_1$ en la
        curva torque--velocidad, y cuál es la no linealidad más evidente de la
        etapa eléctrica?
\end{enumerate}
\end{lapizypapel}

% =====================================================================
\chapter{Dinámica rotacional: dos cuerpos sin eje rígido}\label{cap:mecanica}
% =====================================================================

\begin{ideafisica}
En un motor conectado rígidamente a un ventilador, motor y carga son un solo
bloque sólido que gira siempre al mismo ángulo. En un agitador magnético eso es
falso: abajo gira el rotor con un imán impulsor, y arriba, separado por el fondo
de vidrio del vaso, gira una cápsula magnética libre. \textbf{Son dos cuerpos
independientes}, y cada uno obedece por separado la segunda ley de Newton.
\end{ideafisica}

Cuando uno intenta modelar un sistema mecánico sin dibujar los diagramas de
cuerpo libre por separado, es facilísimo equivocarse en los signos o olvidar una
reacción. Vamos a aplicar el método más seguro que existe en mecánica: aislar
cada cuerpo rígido, dibujar todas las flechas de torque que actúan sobre él y
escribir $\sum T = J\,\dot\omega$.

% ---------------------------------------------------------------------
\section{La segunda ley de Newton en rotación}
% ---------------------------------------------------------------------

En movimiento rectilíneo, la segunda ley de Newton dice $\sum F = m\,\dot v$: la
masa $m$ mide cuánto se resiste un cuerpo a cambiar su velocidad lineal. En
rotación alrededor de un eje fijo, la ley análoga es exactamente igual:
\begin{equation}
  \sum T = J\,\frac{\dd\omega}{\dd t} = J\,\dot\omega,
  \label{eq:newton_rot}
\end{equation}
donde:
\begin{itemize}
  \item $\omega = \dot\theta$ es la velocidad angular [\si{rad/s}] y
        $\dot\omega = \ddot\theta$ es la aceleración angular [\si{rad/s^2}];
  \item $\sum T$ es la suma algebraica de todos los torques externos aplicados
        al cuerpo [\si{N.m}], positivos si empujan en el sentido de giro elegido
        y negativos si lo frenan;
  \item $J$ es el \textbf{momento de inercia} respecto del eje de giro
        [\si{kg.m^2}]. Físicamente, cada trocito de masa $\dd m$ a distancia $r$
        del eje aporta $r^2\,\dd m$ a la inercia ($J = \int r^2\,\dd m$): cuanto
        más masa hay y más lejos del centro está distribuida, más torque cuesta
        acelerarla. Por ejemplo, una barra delgada de masa $m$ y longitud $L$
        girando por su centro tiene $J = \frac{1}{12}m L^2$, y un cilindro
        macizo de radio $r$ tiene $J = \frac{1}{2}m r^2$.
\end{itemize}

\begin{unidades}
Comprobemos que $\sum T = J\,\dot\omega$ es dimensionalmente consistente:
\[
  [J\,\dot\omega] = (\si{kg.m^2})\cdot\left(\frac{\si{rad}}{\si{s^2}}\right)
                  = \left(\si{kg.\frac{m}{s^2}}\right)\cdot\si{m}
                  = \si{N.m} = [T].
\]
\end{unidades}

% ---------------------------------------------------------------------
\section{Diagramas de cuerpo libre y principio de acción y reacción}
% ---------------------------------------------------------------------

Separamos nuestro sistema mecánico en sus dos cuerpos rotantes (Figura~\ref{fig:dcl_dos_inercias}):
\begin{enumerate}
  \item \textbf{Cuerpo~1 (lado motor):} conjunto rígido formado por el rotor del
        BLDC, el soporte y el imán impulsor inferior. Su inercia total es $J_1$,
        su ángulo mecánico es $\theta_1$ y su velocidad angular es $\omega_1 = \dot\theta_1$.
  \item \textbf{Cuerpo~2 (lado agitador):} la cápsula o píldora magnética que
        gira dentro del vaso. Su inercia es $J_2$, su ángulo es $\theta_2$ y su
        velocidad angular es $\omega_2 = \dot\theta_2$.
\end{enumerate}

Definimos la convención de giro positivo (por ejemplo, sentido antihorario visto
desde arriba) y preguntamos qué torques siente cada cuerpo:

\begin{itemize}
  \item \textbf{Sobre el Cuerpo~1 (motor + imán inferior):}
        \begin{itemize}
          \item el campo del estator lo impulsa hacia adelante con el torque
                eléctrico $+T_e$;
          \item sus propios rodamientos lo frenan con la fricción mecánica
                $-T_{f1}(\omega_1)$;
          \item la cápsula de arriba, al ir ligeramente atrasada respecto del
                imán inferior, lo \emph{tira hacia atrás} a través del acople
                magnético con un torque $-T_c$.
        \end{itemize}
  \item \textbf{Sobre el Cuerpo~2 (cápsula agitadora):}
        \begin{itemize}
          \item el imán de abajo la arrastra hacia adelante a través del acople
                con el torque $+T_c$ (\textbf{tercera ley de Newton}: igual en
                magnitud y opuesto en signo al que frena al rotor);
          \item el fluido del vaso se opone a su movimiento con el torque de
                carga $-T_L(\omega_2)$;
          \item el rozamiento de la cápsula contra el fondo del vaso la frena
                con $-T_{f2}(\omega_2)$.
        \end{itemize}
\end{itemize}

\begin{figure}[htbp]
\centering
\begin{tikzpicture}[
  >={Stealth[length=5pt]},
  thick
]
  % ---- Cuerpo 1: Rotor + imán ----
  \begin{scope}[shift={(0,0)}]
    \draw[fill=slatebg, draw=inkblue, very thick] (0,0) ellipse (1.45cm and 0.50cm);
    \node[inkblue, font=\small\bfseries] at (0,0) {Cuerpo 1 ($J_1$)};
    \node[below=0.68cm, font=\footnotesize, align=center] at (0,0) {Rotor + imán impulsor\\$(\theta_1,\,\omega_1)$};

    % Eje vertical
    \draw[dashed, mutedgray] (0,-0.6) -- (0,1.25);

    % Flecha +T_e
    \draw[->, tealaccent, very thick] (-1.65,-0.15) to[bend left=35] (-0.85,0.72);
    \node[above left, tealaccent, font=\small] at (-1.35,0.45) {$+T_e$};

    % Flechas opositoras -T_c y -T_f1
    \draw[->, crimson, very thick] (0.85,0.72) to[bend left=35] (1.65,-0.15);
    \node[above, crimson, font=\scriptsize, align=center] at (1.30,0.68) {$-T_c(\delta)$\\$-T_{f1}(\omega_1)$};
  \end{scope}

  % ---- Conexión acción-reacción ----
  \draw[<->, amberwarn, very thick, decorate, decoration={snake, amplitude=1.6pt, segment length=7pt}]
    (2.35,0.0) -- (6.05,0.0);
  \node[above=0.22cm, amberwarn, font=\footnotesize\bfseries, align=center] at (4.20,0.08) {Acople magnético\\$\delta = \theta_1 - \theta_2$};
  \node[below=0.20cm, amberwarn, font=\scriptsize, align=center] at (4.20,-0.06) {Acción y reacción:\\$-T_c$ en 1, $+T_c$ en 2};

  % ---- Cuerpo 2: Cápsula agitadora ----
  \begin{scope}[shift={(8.4,0)}]
    \draw[fill=amberbg, draw=amberwarn, very thick] (0,0) ellipse (1.40cm and 0.48cm);
    \node[amberwarn, font=\small\bfseries] at (0,0) {Cuerpo 2 ($J_2$)};
    \node[below=0.68cm, font=\footnotesize, align=center] at (0,0) {Cápsula agitadora\\$(\theta_2,\,\omega_2)$};

    % Eje vertical
    \draw[dashed, mutedgray] (0,-0.6) -- (0,1.25);

    % Flecha +T_c
    \draw[->, amberwarn, very thick] (-1.60,-0.15) to[bend left=35] (-0.80,0.70);
    \node[above, amberwarn, font=\scriptsize\bfseries] at (-1.20,0.72) {$+T_c(\delta)$};

    % Flechas opositoras -T_L y -T_f2
    \draw[->, crimson, very thick] (0.80,0.70) to[bend left=35] (1.60,-0.15);
    \node[above right, crimson, font=\scriptsize, align=left] at (1.28,0.40) {$-T_L(\omega_2)$\\$-T_{f2}(\omega_2)$};
  \end{scope}
\end{tikzpicture}
\caption{Diagramas de cuerpo libre separados del rotor ($J_1$) y la cápsula ($J_2$). Aislar cada cuerpo garantiza no equivocarse con el signo de acción y reacción del acople $T_c$.}
\label{fig:dcl_dos_inercias}
\end{figure}

Escribimos ahora $\sum T = J\dot\omega$ para cada cuerpo. Además, lo único que le
importa al acople magnético no es el ángulo absoluto $\theta_1$ ni $\theta_2$
(que crecen sin límite al dar miles de vueltas), sino cuánto se adelantó el
imán inferior respecto de la cápsula. Definimos entonces el \textbf{desfase
angular} $\delta = \theta_1 - \theta_2$, cuya derivada temporal es simplemente
la diferencia de velocidades $\dot\delta = \dot\theta_1 - \dot\theta_2 = \omega_1 - \omega_2$.

Reuniendo las tres ecuaciones mecánicas obtenemos el esqueleto rotacional de
todo el sistema:
\begin{ecuacionclave}
\begin{align}
  J_1\,\dot\omega_1 &= T_e - T_c(\delta) - T_{f1}(\omega_1), \label{eq:newton1}\\
  J_2\,\dot\omega_2 &= T_c(\delta) - T_L(\omega_2) - T_{f2}(\omega_2), \label{eq:newton2}\\
  \dot\delta        &= \omega_1 - \omega_2. \label{eq:cinematica_delta}
\end{align}
\end{ecuacionclave}

% ---------------------------------------------------------------------
\section{Comprobación energética y límite de eje rígido}
% ---------------------------------------------------------------------

Dos comprobaciones rápidas nos confirman que los signos
de~\eqref{eq:newton1}--\eqref{eq:cinematica_delta} son impecables:

\paragraph{1. ¿Qué pasa si los dos cuerpos giran trabados como un solo bloque?}
Si el acople fuese infinitamente rígido ($\delta \equiv 0$, $\omega_1 = \omega_2 = \omega$),
basta sumar~\eqref{eq:newton1} y~\eqref{eq:newton2}: el torque interno $T_c$
aparece con signo menos en una y signo más en la otra, así que \textbf{se cancela
exactamente}:
\[
  (J_1 + J_2)\,\dot\omega = T_e - T_{f1}(\omega) - T_{f2}(\omega) - T_L(\omega).
\]
Recuperamos de inmediato la ecuación de un solo cuerpo con inercia total
$J_{\text{tot}} = J_1 + J_2$. Si hubiéramos puesto mal el signo de $T_c$, no se
habría cancelado.

\paragraph{2. Balance de energía mecánica.}
Multipliquemos~\eqref{eq:newton1} por $\omega_1$, \eqref{eq:newton2} por
$\omega_2$ y sumemos ambas ecuaciones. Recordando que
$J\omega\dot\omega = \frac{\dd}{\dd t}(\tfrac{1}{2}J\omega^2)$ y que
$\omega_1 - \omega_2 = \dot\delta$:
\[
  \frac{\dd}{\dd t}\!\left(\frac{1}{2}J_1\omega_1^2 + \frac{1}{2}J_2\omega_2^2\right)
  + \underbrace{T_c(\delta)\,\dot\delta}_{\dd U_c / \dd t}
  = \underbrace{T_e\,\omega_1}_{\text{potencia entrante}}
  - \underbrace{\big(T_{f1}\omega_1 + T_{f2}\omega_2 + T_L\omega_2\big)}_{\text{potencia disipada}}.
\]
El término $T_c(\delta)\dot\delta$ es la tasa de cambio de la energía potencial
$U_c(\delta) = \int_0^\delta T_c(\alpha)\,\dd\alpha$ almacenada al «retorcer» el
acople. Igual que en el circuito eléctrico, el acople no inventa ni pierde
energía: la guarda cuando se estira y la devuelve cuando se relaja.

\begin{antidoto}
En este capítulo modelamos exclusivamente la rotación en torno al eje vertical.
Dejamos afuera el cabeceo (\emph{wobble}), el desplazamiento radial de la
cápsula hacia las paredes del vaso y los rebotes por excentricidad: son
fenómenos tridimensionales de contacto que no agregan claridad al problema de
control de velocidad y no pueden parametrizarse sin instrumentación óptica o
dinámica de sólidos 3D.
\end{antidoto}

\begin{lapizypapel}
\begin{enumerate}[leftmargin=1.5em]
  \item Dibujá los diagramas de cuerpo libre de $J_1$ y $J_2$ por separado y
        escribí las tres ecuaciones~\eqref{eq:newton1}--\eqref{eq:cinematica_delta}.
  \item ¿Por qué conviene usar $\delta = \theta_1 - \theta_2$ como variable de
        estado en lugar de integrar $\theta_1$ y $\theta_2$ y restarlos después?
  \item Comprobá sumando ambas ecuaciones que cuando $\omega_1 = \omega_2$ el
        torque interno $T_c$ desaparece y las inercias se suman.
\end{enumerate}
\end{lapizypapel}

% =====================================================================
\chapter{El acople magnético: un resorte no lineal que se suelta}\label{cap:acople}
% =====================================================================

\begin{ideafisica}
No hace falta resolver ecuaciones de Maxwell en tres dimensiones para modelar un
acople magnético. Pensalo como un \textbf{problema mecánico y cinemático}: un
resorte torsional que transmite cada vez más torque a medida que lo torcés,
\emph{pero sólo hasta un ángulo crítico}. Si lo torcés más allá de ese ángulo,
los polos se desencuentran, el torque cae y la cápsula se desengancha.
\end{ideafisica}

Hasta ahora dejamos el símbolo $T_c(\delta)$ sin fórmula explícita. En este
capítulo vamos a darle forma a partir de un experimento mental sumamente simple
y descubrir una de las no linealidades más ricas e importantes del sistema: el
\textbf{desacople} (o desincronización) de la cápsula agitadora.

% ---------------------------------------------------------------------
\section{El experimento mental: girar un imán mientras trabás el otro}
% ---------------------------------------------------------------------

Imaginá que sacás el vaso, dejás fija con los dedos la cápsula de arriba
($\theta_2 = 0$) y empezás a girar muy despacio con la otra mano el imán de
abajo ($\theta_1 = \delta$), sintiendo el torque $T_c(\delta)$ con el que el
imán superior intenta seguirte. Para una cápsula magnética típica (un imán
recto con un polo Norte y un polo Sur, es decir $N_a = 1$ par de polos):
\begin{itemize}
  \item \textbf{En $\delta = \SI{0}{\degree}$:} Norte queda enfrentado a Sur y
        Sur a Norte. Es la posición de mínimo reposo: no sentís ningún torque
        ($T_c = 0$).
  \item \textbf{Entre $\SI{0}{\degree}$ y $\SI{90}{\degree}$:} a medida que
        desfasás los imanes, la atracción entre polos opuestos tira con fuerza
        creciente intentando alinearlos de nuevo. El torque crece hasta alcanzar
        un valor pico $T_{\max}$ justo en el \textbf{ángulo crítico}
        $\delta_{\text{crit}} = \SI{90}{\degree}$ ($\pi/2\;\si{rad}$).
  \item \textbf{Entre $\SI{90}{\degree}$ y $\SI{180}{\degree}$:} si seguís
        girando pasado los \SI{90}{\degree}, los polos iguales empiezan a
        acercarse entre sí y la componente de arrastre \emph{disminuye}, hasta
        anularse nuevamente en $\delta = \SI{180}{\degree}$ (Norte frente a
        Norte: equilibrio completamente inestable).
  \item \textbf{Pasados los \SI{180}{\degree}:} el torque cambia de signo: ahora
        el imán superior prefiere «saltar» hacia adelante para alinearse con el
        polo opuesto siguiente ($\SI{360}{\degree}$). El acoplamiento se perdió.
\end{itemize}

Si el acople tuviera $N_a$ pares de polos equiespaciados (por ejemplo, 4 imanes
alternados N--S--N--S, es decir $N_a = 2$), todo este ciclo se repite $N_a$
veces por vuelta mecánica, de modo que el ángulo crítico donde ocurre el pico de
torque es:
\begin{equation}
  \delta_{\text{crit}} = \frac{\pi}{2 N_a}
  \qquad
  \begin{cases}
    N_a = 1 \;\Longrightarrow\; \delta_{\text{crit}} = \SI{90}{\degree},\\
    N_a = 2 \;\Longrightarrow\; \delta_{\text{crit}} = \SI{45}{\degree}.
  \end{cases}
  \label{eq:delta_crit}
\end{equation}

¿Cuál es la función matemática más simple, suave y periódica que vale $0$ en
$\delta=0$, sube monótonamente hasta un máximo $T_{\max}$ en
$\delta_{\text{crit}} = \pi/(2N_a)$, vuelve a $0$ en $2\delta_{\text{crit}}$ y
cambia de signo al cruzarlo? El seno:
\begin{ecuacionclave}
\begin{equation}
  T_c(\delta) = T_{\max}\,\sin(N_a\,\delta)
              = T_{\max}\,\sin\!\left(\frac{\pi}{2}\,\frac{\delta}{\delta_{\text{crit}}}\right).
  \label{eq:acople_seno}
\end{equation}
\end{ecuacionclave}

\paragraph{¿Por qué justamente un seno y no otra curva parecida?}
Además de ser la ley cinemática más limpia que describe el experimento mental,
tiene un respaldo físico elemental: la energía potencial de un dipolo magnético
en un campo uniforme es $U_c(\delta) = -U_1\cos(N_a\delta)$ (mínima en
$\delta=0$, máxima en $\delta=\pi/N_a$). Por conservación de energía, el torque
necesario para sostener un desfase $\delta$ es la derivada de esa energía
respecto del ángulo:
\[
  T_c(\delta) = \frac{\dd U_c}{\dd\delta} = N_a U_1\sin(N_a\delta) \equiv T_{\max}\sin(N_a\delta).
\]
No hace falta calcular $U_1$ a partir de integrales volumétricas de campo:
$T_{\max}$ (el \emph{torque de desenganche} o \emph{pull-out torque}) es un
único número físico medible en un banco (por ejemplo, midiendo qué fuerza en un
brazo de palanca despega el imán) y $N_a$ se conoce mirando la geometría del
imán. Dos parámetros concretos, cero misterios.

% ---------------------------------------------------------------------
\section{Equilibrio estático, estabilidad y desenganche}
% ---------------------------------------------------------------------

Supongamos que motor y cápsula giran sincronizados a una velocidad constante
$\omega_1 = \omega_2 = \omega_0$. Entonces $\dot\omega_2 = 0$ y $\dot\delta = 0$.
Llamemos $T_{L,\text{tot}}(\omega_0) = T_L(\omega_0) + T_{f2}(\omega_0)$ a toda
la carga que frena a la cápsula. De la ecuación~\eqref{eq:newton2} del Cuerpo~2
resulta la condición de equilibrio:
\begin{equation}
  T_{\max}\,\sin(N_a\,\delta^\ast) = T_{L,\text{tot}}(\omega_0).
  \label{eq:eq_acople_cap3}
\end{equation}

Miremos la Figura~\ref{fig:curva_acople}: resolver~\eqref{eq:eq_acople_cap3}
equivale a buscar dónde la recta horizontal de carga $T_{L,\text{tot}}$ corta a
la curva senoidal $T_c(\delta)$. Aparecen inmediatamente dos regímenes
radicalmente distintos:

\begin{figure}[htbp]
\centering
\begin{tikzpicture}
\begin{axis}[
  width=11.6cm, height=6.6cm,
  axis lines=middle,
  xlabel={Desfase angular $N_a\delta$},
  xlabel style={at={(ticklabel* cs:1.0)}, anchor=north east, yshift=-18pt},
  ylabel={$T_c(\delta)$},
  ylabel style={at={(ticklabel* cs:1.02)}, anchor=south},
  xmin=-0.3, xmax=3.65, ymin=-0.45, ymax=1.38,
  xtick={0, 0.5236, 1.5708, 2.6180, 3.1416},
  xticklabels={$0$, $N_a\delta^\ast$, $\frac{\pi}{2}$ ($\delta_{\text{crit}}$), $N_a\delta_u$, $\pi$},
  ytick={0.5, 1.0},
  yticklabels={$T_{L,\text{tot}}$, $T_{\max}$},
  tick label style={font=\footnotesize},
  label style={font=\small},
  clip=false
]
  % Zona estable sombreada
  \fill[tealaccent!8] (axis cs:0,0) rectangle (axis cs:1.5708,1.08);
  \fill[crimson!7] (axis cs:1.5708,0) rectangle (axis cs:3.1416,1.08);

  \node[tealaccent, font=\scriptsize\bfseries, align=center] at (axis cs:0.82,1.22) {Rama estable\\(resorte rígido $\to$ blando)};
  \node[crimson, font=\scriptsize\bfseries, align=center] at (axis cs:2.38,1.22) {Rama inestable\\(desenganche)};

  % Curva senoidal T_max sin(x)
  \addplot[inkblue, very thick, domain=-0.25:3.5, samples=120] {sin(deg(x))};

  % Techo T_max
  \draw[crimson, dashed, thin] (axis cs:0,1.0) -- (axis cs:3.4,1.0);
  \draw[mutedgray, dotted] (axis cs:1.5708,0) -- (axis cs:1.5708,1.0);

  % Línea de carga T_L,tot = 0.5
  \draw[amberwarn, thick, dashed] (axis cs:0,0.5) -- (axis cs:3.3,0.5);

  % Tangente linealizada en delta* = pi/6 (cos(pi/6)=0.866)
  \addplot[tealaccent, thick, domain=0.05:1.10] {0.5 + 0.866*(x - 0.5236)};
  \node[tealaccent, font=\scriptsize, left] at (axis cs:0.88,0.88) {pendiente $K_{\text{eff}}$};

  % Puntos de equilibrio
  \fill[tealaccent] (axis cs:0.5236,0.5) circle (3pt);
  \draw[tealaccent, dotted] (axis cs:0.5236,0) -- (axis cs:0.5236,0.5);
  \node[below right, tealaccent, font=\scriptsize\bfseries] at (axis cs:0.56,0.46) {Equil. estable};

  \fill[crimson] (axis cs:2.6180,0.5) circle (3pt);
  \draw[crimson, dotted] (axis cs:2.6180,0) -- (axis cs:2.6180,0.5);
  \node[above right, crimson, font=\scriptsize\bfseries] at (axis cs:2.65,0.54) {Equil. inestable};
\end{axis}
\end{tikzpicture}
\caption{Curva torque--desfase del acople magnético $T_c(\delta) = T_{\max}\sin(N_a\delta)$. Para cualquier carga $T_{L,\text{tot}} < T_{\max}$ existen dos intersecciones: sólo la primera ($\delta^\ast < \delta_{\text{crit}}$) es estable. La recta tangente verde muestra la rigidez linealizada $K_{\text{eff}}$.}
\label{fig:curva_acople}
\end{figure}

\paragraph{Caso 1: Carga admisible ($T_{L,\text{tot}} < T_{\max}$).}
La recta horizontal corta al seno en \textbf{dos puntos}:
\[
  \delta^\ast = \frac{1}{N_a}\arcsin\!\left(\frac{T_{L,\text{tot}}}{T_{\max}}\right) < \delta_{\text{crit}},
  \qquad\text{y}\qquad
  \delta_u = \frac{\pi}{N_a} - \delta^\ast > \delta_{\text{crit}}.
\]
¿Por qué sólo $\delta^\ast$ es estable y $\delta_u$ no lo es? Hagamos la prueba
física más sencilla del mundo: \emph{empuje pequeño} $\delta = \delta_{\text{eq}} + \tilde\delta$
(con $\tilde\delta > 0$, es decir, la cápsula se atrasó un poquito más):
\begin{itemize}
  \item \textbf{Alrededor de $\delta^\ast$ (pendiente positiva):} al aumentar
        $\delta$, el torque del acople \emph{aumenta} ($T_c > T_{L,\text{tot}}$).
        Ese exceso de torque acelera la cápsula ($\dot\omega_2 > 0$), haciendo
        que alcance al motor y reduzca $\delta$ de vuelta hacia $\delta^\ast$.
        \textbf{El equilibrio se corrige solo: es estable.}
  \item \textbf{Alrededor de $\delta_u$ (pendiente negativa):} al aumentar
        $\delta$, el torque del acople \emph{disminuye} ($T_c < T_{L,\text{tot}}$).
        Como ahora la carga gana al acople, la cápsula se frena aún más
        ($\dot\omega_2 < 0$), lo que hace crecer todavía más $\delta$.
        \textbf{El error se amplifica solo: es inestable.}
\end{itemize}

\paragraph{Caso 2: Sobrecarga ($T_{L,\text{tot}} > T_{\max}$).}
La recta horizontal queda por encima del pico $T_{\max}$: \textbf{no existe
ninguna intersección}. No hay ningún ángulo capaz de sostener la carga.
Entonces $\omega_2$ cae por debajo de $\omega_1$, el desfase $\delta(t)$ crece
sin cota y $T_c(\delta) = T_{\max}\sin(N_a\delta)$ pasa a oscilar rápidamente
entre $+T_{\max}$ y $-T_{\max}$ con \textbf{valor medio temporal nulo}. Al no
recibir torque medio neto, la cápsula se frena hasta detenerse ($\omega_2 \to 0$)
o vibrar erráticamente en el fondo del vaso, aunque el motor siga girando rápido
abajo.

% ---------------------------------------------------------------------
\section{Rigidez local y «ablandamiento» del resorte magnético}
% ---------------------------------------------------------------------

Si el sistema trabaja cerca de un equilibrio estable $\delta^\ast$ y sufre
pequeñas perturbaciones $\tilde\delta = \delta - \delta^\ast$, la curva
senoidal se parece localmente a su recta tangente (Figura~\ref{fig:curva_acople}).
Su pendiente local es la \textbf{rigidez torsional equivalente}:
\begin{equation}
  K_{\text{eff}}(\delta^\ast)
  = \left.\frac{\dd T_c}{\dd\delta}\right|_{\delta^\ast}
  = N_a\,T_{\max}\,\cos(N_a\,\delta^\ast).
  \label{eq:keff_cos}
\end{equation}
Usando la identidad $\cos\alpha = \sqrt{1 - \sin^2\alpha}$ y sustituyendo
$\sin(N_a\delta^\ast) = T_{L,\text{tot}}/T_{\max}$ de~\eqref{eq:eq_acople_cap3},
la podemos expresar directamente en función de la carga:
\begin{ecuacionclave}
\begin{equation}
  K_{\text{eff}}(\delta^\ast)
  = N_a\,T_{\max}\,\cos(N_a\,\delta^\ast)
  = N_a\,\sqrt{T_{\max}^2 - T_{L,\text{tot}}^2}.
  \label{eq:keff_carga}
\end{equation}
\end{ecuacionclave}

\begin{lectura}
Mirá qué enseñanza física tan fuerte da~\eqref{eq:keff_carga}:
\begin{itemize}
  \item \textbf{En vacío ($T_{L,\text{tot}} = 0$, $\delta^\ast = 0$):} el acople
        está en su punto más rígido, $K_{\text{eff}} = N_a T_{\max}$.
  \item \textbf{Al aumentar la carga ($T_{L,\text{tot}} \to T_{\max}$):} la
        raíz cuadrada se achica y el «resorte magnético» se vuelve cada vez
        \textbf{más blando}, hasta que en $\delta^\ast = \delta_{\text{crit}}$
        su rigidez cae a cero ($K_{\text{eff}} = 0$).
\end{itemize}
Por eso, linealizar el acople en $\delta = 0$ sólo sirve en vacío; con carga en
el vaso, la dinámica torsional cambia con el punto de operación.
\end{lectura}

\paragraph{Frecuencia natural de oscilación torsional.}
Cuando dos inercias $J_1$ y $J_2$ están unidas por un resorte de rigidez
$K_{\text{eff}}$, pueden oscilar una contra la otra (como dos volantes unidos
por una barra elástica). Para hallar esa frecuencia natural, perturbemos
ligeramente el equilibrio dejando de lado por un instante los torques externos
y las fricciones:
\[
  J_1\,\dot{\tilde\omega}_1 = -K_{\text{eff}}\,\tilde\delta,
  \qquad
  J_2\,\dot{\tilde\omega}_2 = +K_{\text{eff}}\,\tilde\delta.
\]
Derivando $\dot{\tilde\delta} = \tilde\omega_1 - \tilde\omega_2$ respecto del
tiempo y sustituyendo ambas aceleraciones:
\[
  \ddot{\tilde\delta}
  = \dot{\tilde\omega}_1 - \dot{\tilde\omega}_2
  = -\frac{K_{\text{eff}}}{J_1}\tilde\delta - \frac{K_{\text{eff}}}{J_2}\tilde\delta
  = -K_{\text{eff}}\left(\frac{1}{J_1} + \frac{1}{J_2}\right)\tilde\delta
  = -\frac{K_{\text{eff}}}{J_{\text{eq}}}\,\tilde\delta,
\]
donde definimos la \textbf{inercia reducida} (equivalente serie):
\begin{equation}
  J_{\text{eq}} = \frac{J_1\,J_2}{J_1 + J_2},
  \qquad\Longrightarrow\qquad
  \omega_{n,\text{tors}}(\delta^\ast) = \sqrt{\frac{K_{\text{eff}}(\delta^\ast)}{J_{\text{eq}}}}.
  \label{eq:wn_torsional}
\end{equation}
Verificación de sentido común: si el rotor es mucho más pesado que la cápsula
($J_1 \gg J_2$), entonces $J_{\text{eq}} \approx J_2$: la cápsula liviana vibra
contra un rotor casi inamovible a frecuencia $\sqrt{K_{\text{eff}}/J_2}$. Y a
medida que la carga acerca $\delta^\ast$ a $\delta_{\text{crit}}$, $K_{\text{eff}}$
baja y esa vibración torsional se vuelve más lenta y floja antes de soltarse.

% ---------------------------------------------------------------------
\section{Desacople dinámico: por qué un acelerón suelta la barra aunque \texorpdfstring{$T_{L,\text{tot}} < T_{\max}$}{TL,tot < Tmax}}
% ---------------------------------------------------------------------

Hay un fenómeno que cualquiera que haya usado un agitador magnético vio en el
laboratorio: la carga del líquido está bien por debajo de $T_{\max}$ (en régimen
permanente el agitador podría girar sin problema), pero si le pegás un salto
brusco de velocidad al motor, \textbf{la cápsula se suelta igual}. ¿Por qué?

Podemos explicarlo en cinco líneas con lápiz y papel. Supongamos que el motor
da un salto repentino de velocidad y se queda girando a una velocidad constante
$\omega_1$ ($\dot\omega_1 = 0$), mientras la cápsula todavía viene más lenta
($\dot\delta = \omega_1 - \omega_2 > 0$). Tomando $T_{L,\text{tot}}$ casi
constante durante ese breve instante, como $\ddot\delta = \dot\omega_1 - \dot\omega_2 = -\dot\omega_2$,
multiplicamos la ecuación~\eqref{eq:newton2} por $-1$:
\begin{equation}
  J_2\,\ddot\delta = -T_{\max}\,\sin(N_a\,\delta) + T_{L,\text{tot}}.
  \label{eq:pendulo_acople}
\end{equation}
¡Es matemáticamente idéntica a la ecuación de un \textbf{péndulo} empujado por
un torque constante! Para ver cuándo la cápsula «da la vuelta» (se desengancha),
buscamos su conservación de energía multiplicando ambos lados por $\dot\delta$:
\[
  J_2\,\dot\delta\,\ddot\delta + T_{\max}\sin(N_a\delta)\,\dot\delta - T_{L,\text{tot}}\,\dot\delta = 0.
\]
Cada término es una derivada exacta respecto del tiempo:
\[
  \frac{\dd}{\dd t}\!\left[
    \underbrace{\frac{1}{2}\,J_2\,\dot\delta^2}_{\text{energía cinética relativa}}
    + \underbrace{\left(-\frac{T_{\max}}{N_a}\cos(N_a\delta) - T_{L,\text{tot}}\,\delta\right)}_{U_{\text{eff}}(\delta)}
  \right] = 0.
\]
Por lo tanto, durante el transitorio se conserva la energía total
$E = \frac{1}{2}J_2\dot\delta^2 + U_{\text{eff}}(\delta)$.

La Figura~\ref{fig:pozo_potencial} dibuja esa energía potencial efectiva
$U_{\text{eff}}(\delta)$: es un \textbf{pozo de potencial inclinado} (como una
chapa ondulada puesta en bajada). El fondo del valle está justo en el equilibrio
estable $\delta^\ast$, y la cima de la colina siguiente está justo en el
equilibrio inestable $\delta_u$:
\begin{itemize}
  \item Si el salto de velocidad aporta una energía cinética relativa
        $\frac{1}{2}J_2\dot\delta^2$ \textbf{menor} que la altura de la barrera
        $\Delta U = U_{\text{eff}}(\delta_u) - U_{\text{eff}}(\delta^\ast)$, la
        cápsula oscila dentro del pozo y vuelve a estabilizarse en $\delta^\ast$.
  \item Si aporta energía \textbf{mayor} que $\Delta U$, la bolita pasa por
        encima de la colina $\delta_u$ y cae indefinidamente cuesta abajo:
        \textbf{la cápsula se desacopla}.
\end{itemize}

\begin{figure}[htbp]
\centering
\begin{tikzpicture}
\begin{axis}[
  width=11.5cm, height=5.8cm,
  axis lines=left,
  xlabel={Desfase angular $N_a\delta$},
  ylabel={Potencial efectivo $U_{\text{eff}}(\delta)$ (norm.)},
  xmin=-0.6, xmax=5.2, ymin=-2.8, ymax=0.9,
  xtick={0.5236, 2.6180},
  xticklabels={$N_a\delta^\ast$ (valle), $N_a\delta_u$ (barrera)},
  ytick=\empty,
  tick label style={font=\footnotesize},
  label style={font=\small},
  clip=false
]
  % U_eff(x) = -cos(x) - 0.5*x (para T_L/T_max = 0.5)
  \addplot[inkblue, very thick, domain=-0.5:5.0, samples=140] {-cos(deg(x)) - 0.5*x};

  % Niveles del valle y la barrera
  % En x* = pi/6 = 0.5236: -cos(30) - 0.2618 = -0.866 - 0.262 = -1.128
  % En xu = 5pi/6 = 2.6180: -cos(150) - 1.309 = +0.866 - 1.309 = -0.443
  \draw[tealaccent, dashed] (axis cs:-0.2,-1.128) -- (axis cs:1.8,-1.128);
  \draw[crimson, dashed] (axis cs:0.5236,-0.443) -- (axis cs:3.4,-0.443);

  % Altura de barrera Delta U
  \draw[<->, amberwarn, very thick] (axis cs:0.5236,-1.128) -- (axis cs:0.5236,-0.443)
    node[midway, right, font=\footnotesize\bfseries]{$\Delta U$};

  % Puntos
  \fill[tealaccent] (axis cs:0.5236,-1.128) circle (3pt);
  \fill[crimson] (axis cs:2.6180,-0.443) circle (3pt);

  % Trayectorias conceptuales
  \node[tealaccent, font=\scriptsize, align=center] at (axis cs:1.45,-1.65) {Energía $< \Delta U$:\\oscila y queda enganchada};
  \draw[->, crimson, thick] (axis cs:2.85,-0.35) to[bend left=18] (axis cs:4.1,-1.55);
  \node[crimson, font=\scriptsize, align=left] at (axis cs:4.15,-0.25) {Energía $> \Delta U$:\\escapa y desacopla!};
\end{axis}
\end{tikzpicture}
\caption{Energía potencial efectiva del acople bajo carga constante ($T_{L,\text{tot}} = 0{,}5\,T_{\max}$). Cuanto mayor es la carga, más inclinada queda la curva y más baja se vuelve la barrera $\Delta U$ que impide el desenganche.}
\label{fig:pozo_potencial}
\end{figure}

Incluso en vacío ($T_{L,\text{tot}} = 0$), donde $\delta^\ast = 0$ y
$\delta_u = \pi/N_a$, la barrera mide $\Delta U = \frac{2T_{\max}}{N_a}$. Por
ende, el mayor salto instantáneo de velocidad que la cápsula puede soportar sin
soltarse es:
\[
  \frac{1}{2}J_2\,(\Delta\omega_{\max})^2 = \frac{2T_{\max}}{N_a}
  \quad\Longrightarrow\quad
  \Delta\omega_{\max} = 2\sqrt{\frac{T_{\max}}{N_a\,J_2}}.
\]
Con esta sola línea se entiende por qué una \textbf{rampa suave} de arranque
mantiene siempre a la cápsula cerca del fondo del valle, mientras que un escalón
brusco la expulsa por encima de la barrera.

\begin{antidoto}
Dejamos afuera intencionalmente:
\begin{itemize}
  \item \textbf{Cálculo magnetostático 3D de $T_{\max}$ en función de la
        geometría y el entrehierro:} requiere conocer la remanencia exacta $B_r$
        de imanes comerciales sin hoja de datos y resolver integrales de dipolos.
        En cambio, $T_{\max}$ es una única magnitud mecánica medible o acotable
        conservadoramente.
  \item \textbf{Amortiguamiento magnético interno ($c\,\dot\delta$) por
        corrientes parásitas:} en un vaso de vidrio y una cápsula recubierta en
        PTFE es despreciable, y no agregar amortiguamiento no medido vuelve al
        modelo más conservador frente a oscilaciones y desenganche.
\end{itemize}
\end{antidoto}

\begin{lapizypapel}
\begin{enumerate}[leftmargin=1.5em]
  \item Dibujá la curva $T_c(\delta) = T_{\max}\sin(N_a\delta)$ junto a una
        carga $T_{L,\text{tot}} < T_{\max}$. Señalá $\delta_{\text{crit}}$,
        $\delta^\ast$ y $\delta_u$, y justificá físicamente por qué uno es
        estable y el otro inestable.
  \item Deducí la rigidez linealizada $K_{\text{eff}} = N_a\sqrt{T_{\max}^2 - T_{L,\text{tot}}^2}$
        y explicá qué le pasa al «resorte magnético» al acercarse a $T_{\max}$.
  \item Explicá, sin usar fórmulas complicadas, por qué un escalón agresivo de
        velocidad puede desacoplar la cápsula aun cuando $T_{L,\text{tot}} < T_{\max}$.
\end{enumerate}
\end{lapizypapel}

% =====================================================================
\chapter{Las fuerzas que frenan: fricción mecánica y arrastre del fluido}\label{cap:cargas}
% =====================================================================

\begin{ideafisica}
Si no existieran rozamientos ni fluido, un motor al que le cortás la corriente
giraría eternamente. En la realidad dos cosas disipan energía mecánica y exigen
torque permanente para mantener la velocidad: el \textbf{rozamiento sólido}
(rodamientos y contacto de la cápsula con el vaso) y la \textbf{resistencia del
fluido} que la cápsula tiene que empujar y revolver.
\end{ideafisica}

En los diagramas del Capítulo~\ref{cap:mecanica} aparecieron tres torques
resistivos: $T_{f1}(\omega_1)$, $T_{f2}(\omega_2)$ y $T_L(\omega_2)$. Veamos
cómo se modelan de la manera más sencilla, físicamente clara y fácil de
identificar.

% ---------------------------------------------------------------------
\section{Fricción mecánica: seca (Coulomb) más viscosa}
% ---------------------------------------------------------------------

Al mover cualquier pieza mecánica sobre un apoyo o rodamiento aparecen dos
contribuciones clásicas:
\begin{enumerate}
  \item \textbf{Fricción seca o de Coulomb ($T_C$):} es un torque de magnitud
        constante que siempre se opone al sentido del movimiento, sin importar
        qué tan rápido gire el eje ($T_C\,\sgn(\omega)$). Físicamente proviene
        del rozamiento sólido entre superficies (bolillas y pistas del
        rodamiento, o la cápsula apoyada por su propio peso sobre el fondo de
        vidrio del vaso).
  \item \textbf{Fricción viscosa ($b\,\omega$):} es un torque proporcional a la
        velocidad angular, debido al cizallamiento de películas de lubricante o
        aire en el rodamiento y el apoyo.
\end{enumerate}

Combinando ambas para cada cuerpo $i \in \{1, 2\}$:
\begin{ecuacionclave}
\begin{equation}
  T_{f,i}(\omega_i) = T_{C,i}\,\tanh\!\left(\frac{\omega_i}{\omega_{0,f}}\right) + b_i\,\omega_i,
  \qquad i \in \{1, 2\}.
  \label{eq:friccion_cap4}
\end{equation}
\end{ecuacionclave}

\paragraph{¿Por qué escribimos $\tanh(\omega_i/\omega_{0,f})$ en lugar de $\sgn(\omega_i)$?}
No es ningún fenómeno físico oculto ni misterioso: la función signo $\sgn(\omega)$
pega un salto vertical discontinuo en $\omega = 0$ que hace «vibrar» a los
métodos numéricos de integración cuando el cuerpo se frena. La tangente
hiperbólica $\tanh(\omega/\omega_{0,f})$, eligiendo una banda pequeña
$\omega_{0,f}$ frente a las velocidades de trabajo (por ejemplo
$\omega_{0,f} \approx \SI{5}{rad/s}$), vale prácticamente $+1$ para toda
velocidad positiva normal ($\tanh(3) \approx 0{,}995$), vale $-1$ para
velocidad negativa y cruza suavemente por $0$ en reposo. En todo el rango de
operación del motor, $T_{C,i}\tanh(\omega_i/\omega_{0,f}) \approx T_{C,i}$.

\begin{unidades}
Verificamos las unidades de ambos parámetros:
\[
  [T_{C,i}] = \si{N.m}, \qquad
  [b_i] = \frac{[T]}{[\omega]} = \frac{\si{N.m}}{\si{rad/s}} = \si{N.m.s/rad}.
\]
\end{unidades}

% ---------------------------------------------------------------------
\section{Cómo separar \texorpdfstring{$T_C$ y $b$}{TC y b} en la práctica: ensayo de desaceleración libre}
% ---------------------------------------------------------------------

Una gran ventaja de este modelo de dos parámetros es que \textbf{no hay que
adivinarlos}: se ven a simple vista en un ensayo de desaceleración libre
(\emph{coast-down}). Imaginemos que llevamos el motor (en vacío, sin cápsula) a
una velocidad $\omega_0 > 0$ y cortamos por completo la alimentación ($T_e = 0$).
Mientras siga girando ($\omega > 0$), la ecuación de Newton queda:
\begin{equation}
  J\,\frac{\dd\omega}{\dd t} + b\,\omega = -T_C.
  \label{eq:edo_coastdown}
\end{equation}
Resolvamos esta ecuación diferencial lineal de primer orden paso a paso:
\begin{enumerate}
  \item \textbf{Solución homogénea} ($J\dot\omega_h + b\omega_h = 0$):
        $\omega_h(t) = A\,e^{-t/\tau}$, con constante de tiempo $\tau = J/b$.
  \item \textbf{Solución particular constante} ($\dot\omega_p = 0$):
        $b\,\omega_p = -T_C \implies \omega_p = -T_C/b$.
  \item \textbf{Solución general:} $\omega(t) = A\,e^{-t/\tau} - \frac{T_C}{b}$.
  \item \textbf{Condición inicial $\omega(0) = \omega_0$:}
        $A - \frac{T_C}{b} = \omega_0 \implies A = \omega_0 + \frac{T_C}{b}$.
\end{enumerate}
Por lo tanto, mientras $\omega(t) \ge 0$:
\begin{equation}
  \boxed{\;\omega(t) = \left(\omega_0 + \frac{T_C}{b}\right)e^{-t/\tau} - \frac{T_C}{b},
  \qquad \tau = \frac{J}{b}\;.}
  \label{eq:sol_coastdown}
\end{equation}

Mirá los dos casos límite extremos en la Figura~\ref{fig:coastdown_fluido}
(izquierda):
\begin{itemize}
  \item Si domina la \textbf{fricción seca} ($b \to 0$), la desaceleración es
        constante ($\dot\omega = -T_C/J$) y la velocidad cae en \textbf{línea
        recta} hasta detenerse en un tiempo finito $t_{\text{stop}} = J\omega_0/T_C$.
  \item Si domina la \textbf{fricción viscosa} ($T_C \to 0$), la velocidad cae
        como una \textbf{exponencial} $\omega_0 e^{-t/\tau}$ que se frena rápido
        al principio y se achata lentamente hacia cero.
\end{itemize}
La curvatura de la desaceleración libre te dice directamente cuánto pesa cada
uno!

% ---------------------------------------------------------------------
\section{El torque de carga del fluido: viscosidad frente a inercia}
% ---------------------------------------------------------------------

Ahora llenemos el vaso con líquido. ¿Qué siente la cápsula agitadora al girar a
velocidad $\omega_2$? Pensá en revolver con una cuchara dos líquidos bien
distintos: miel espesa despacio, frente a agua rápido. Estás sintiendo dos
mecanismos físicos diferentes:

\begin{enumerate}
  \item \textbf{Arrastre viscoso (régimen laminar):} el líquido se adhiere a la
        superficie de la cápsula y las capas vecinas de fluido se deslizan unas
        sobre otras. Por la ley de viscosidad de Newton, el esfuerzo cortante es
        proporcional a la viscosidad dinámica $\mu$ y a la velocidad de
        deformación (proporcional a $\omega_2$). Esto produce un torque
        \textbf{lineal en la velocidad}:
        \[
          T_{L,\text{visc}}(\omega_2) = b_{fl}\,\omega_2 \qquad (b_{fl} \propto \mu).
        \]
  \item \textbf{Arrastre inercial o de forma (régimen turbulento):} al girar más
        rápido, la cápsula tiene que empujar y acelerar continuamente masa de
        fluido que está en su camino. En un instante $\dd t$, la masa de fluido
        barrida es proporcional a la densidad $\rho$ y a la velocidad lineal
        $v \propto \omega_2$, y esa masa se acelera a una velocidad también
        proporcional a $\omega_2$. Por variación de cantidad de movimiento (o
        presión dinámica $\frac{1}{2}\rho v^2$), la fuerza crece como $\rho v^2$,
        generando un torque \textbf{cuadrático en la velocidad}:
        \[
          T_{L,\text{iner}}(\omega_2) = c_{fl}\,\omega_2\,|\omega_2| \qquad (c_{fl} \propto \rho).
        \]
        \emph{Detalle de signo:} escribimos $\omega_2|\omega_2|$ y no $\omega_2^2$
        porque $\omega_2^2$ sería siempre positivo; en cambio, $\omega_2|\omega_2|$
        cambia de signo si la cápsula gira al revés, garantizando que el fluido
        \textbf{siempre frene} el movimiento y nunca lo empuje.
\end{enumerate}

Sumando ambos efectos obtenemos el modelo de carga fluida:
\begin{ecuacionclave}
\begin{equation}
  T_L(\omega_2) = b_{fl}\,\omega_2 + c_{fl}\,\omega_2\,|\omega_2|.
  \label{eq:carga_fluido}
\end{equation}
\end{ecuacionclave}

\begin{unidades}
Comprobamos las dimensiones de ambos coeficientes efectivos del fluido:
\[
  [b_{fl}] = \si{N.m.s/rad}, \qquad
  [c_{fl}] = \frac{\si{N.m}}{\si{rad^2/s^2}} = \si{N.m.s^2/rad^2} = \si{kg.m^2/rad^2}.
\]
Notá la consistencia física: como $c_{fl} \propto \rho\,D^5$ (densidad por
longitud a la quinta), sus unidades son $(\si{kg/m^3})\cdot\si{m^5} = \si{kg.m^2}$!
\end{unidades}

\begin{figure}[htbp]
\centering
\begin{minipage}[t]{0.48\linewidth}
\centering
\begin{tikzpicture}
\begin{axis}[
  width=6.6cm, height=5.0cm,
  axis lines=left,
  xlabel={Tiempo $t$},
  ylabel={Velocidad $\omega(t)$},
  xmin=0, xmax=3.2, ymin=0, ymax=1.15,
  xtick=\empty, ytick={1}, yticklabels={$\omega_0$},
  tick label style={font=\footnotesize},
  label style={font=\footnotesize},
  legend style={font=\scriptsize, at={(0.98,0.98)}, anchor=north east, draw=none, fill=white, fill opacity=0.85}
]
  % Pura viscosa
  \addplot[tealaccent, thick, dashed, domain=0:3.1, samples=60] {exp(-0.7*x)};
  \addlegendentry{Sólo viscosa ($b$)}
  % Pura Coulomb
  \addplot[crimson, thick, dotted, domain=0:2.0, samples=30] {1 - 0.5*x};
  \addlegendentry{Sólo seca ($T_C$)}
  % Combinada
  \addplot[inkblue, very thick, domain=0:1.45, samples=60] {max(0, 1.5*exp(-0.7*x) - 0.5)};
  \addlegendentry{Combinada ($T_C + b\omega$)}
\end{axis}
\end{tikzpicture}
\end{minipage}
\hfill
\begin{minipage}[t]{0.48\linewidth}
\centering
\begin{tikzpicture}
\begin{axis}[
  width=6.6cm, height=5.0cm,
  axis lines=left,
  xlabel={Velocidad $\omega_2$},
  ylabel={Torque resistivo},
  xmin=0, xmax=1.15, ymin=0, ymax=1.25,
  xtick=\empty, ytick={0.18}, yticklabels={$T_{C2}$},
  tick label style={font=\footnotesize},
  label style={font=\footnotesize},
  legend style={font=\scriptsize, at={(0.03,0.98)}, anchor=north west, draw=none, fill=white, fill opacity=0.85}
]
  % Fricción mecánica Tc + b*w
  \addplot[tealaccent, thick, dashed, domain=0:1.1, samples=60] {0.18*tanh(12*x) + 0.25*x};
  \addlegendentry{Fricción $T_{f2}$}
  % Carga fluido bfl*w + cfl*w^2
  \addplot[amberwarn, thick, domain=0:1.1, samples=60] {0.22*x + 0.55*x^2};
  \addlegendentry{Fluido $T_L$}
  % Total
  \addplot[inkblue, very thick, domain=0:1.1, samples=60] {0.18*tanh(12*x) + 0.47*x + 0.55*x^2};
  \addlegendentry{Total $T_{L,\text{tot}}$}
\end{axis}
\end{tikzpicture}
\end{minipage}
\caption{Izquierda: desaceleración libre (\emph{coast-down}) comparando fricción viscosa, seca y combinada. Derecha: suma de fricción mecánica y arrastre cuadrático del fluido sobre la cápsula.}
\label{fig:coastdown_fluido}
\end{figure}

% ---------------------------------------------------------------------
\section{El escenario crítico: pasar de aire a líquido y la pérdida de observabilidad}
% ---------------------------------------------------------------------

Miremos ahora qué le sucede al sistema completo cuando arrancamos el agitador
en un vaso vacío (aire: $b_{fl}\approx 0$, $c_{fl}\approx 0$) y, estando ya en
régimen a velocidad $\omega_0$, vertemos líquido o aumentamos bruscamente su
viscosidad en un instante $t_0$:
\[
  b_{fl}(t) = b_{fl,0} + \Delta b_{fl}\,H(t - t_0),
  \qquad
  c_{fl}(t) = c_{fl,0} + \Delta c_{fl}\,H(t - t_0),
\]
donde $H(t - t_0)$ es la función escalón unitario ($0$ para $t < t_0$ y $1$ para $t \ge t_0$).

Sigamos las ecuaciones~\eqref{eq:newton1}--\eqref{eq:newton2} línea por línea:
\begin{enumerate}
  \item Al ingresar el líquido en $t = t_0$, la carga $T_L(\omega_2)$ pega un
        salto repentino. La cápsula se frena, $\omega_2$ baja y el desfase
        $\delta = \theta_1 - \theta_2$ empieza a crecer.
  \item Si la nueva carga supera el techo del acople ($T_{L,\text{tot}}(\omega_0) > T_{\max}$),
        o si el salto dinámico alcanza para superar la barrera $\Delta U$ de la
        Figura~\ref{fig:pozo_potencial}, la cápsula se desengancha: $\delta$
        gira libremente, el promedio temporal $\langle T_c(\delta)\rangle \to 0$
        y la cápsula se detiene ($\omega_2 \to 0$).
  \item \textbf{¿Qué ve el motor?} Volvé a mirar la ecuación del Cuerpo~1:
        \[
          J_1\,\dot\omega_1 = T_e - \underbrace{T_c(\delta)}_{\langle T_c\rangle \to 0} - T_{f1}(\omega_1).
        \]
        Al desaparecer el torque de reacción del acople, \textbf{el motor queda
        mecánicamente desconectado de la cápsula}. Un controlador que sólo
        mide la velocidad del rotor $\omega_1$ (sin medir corriente ni ver la
        cápsula) mantiene felizmente $\omega_1 \approx \omega^\ast$ y
        \textbf{jamás se entera de que el agitador dejó de agitar}.
\end{enumerate}

Este fenómeno ---capturable por el modelo aunque invisible para el sensor de
velocidad del motor--- es exactamente una de las conclusiones de ingeniería más
valiosas que el modelado permite demostrar antes de tocar el hardware.

\begin{antidoto}
En este capítulo dejamos afuera deliberadamente:
\begin{itemize}
  \item \textbf{Fricción de Stribeck y modelo dinámico de LuGre:} describen el
        pico de fricción estática y microdesplazamientos cerca de $\omega = 0$
        (\emph{stick-slip}). Como el sensor de velocidad está ciego a bajas rpm
        y el arranque se hace a lazo abierto, esos parámetros adicionales son
        inidentificables con la instrumentación real y no intervienen en régimen.
  \item \textbf{Tablas industriales de número de potencia $N_p(\mathrm{Re})$:}
        esas correlaciones están tabuladas para tanques normalizados con
        deflectores (\emph{baffles}) y turbinas Rushton, no para una cápsula
        magnética en un vaso de precipitados liso (donde además se forma un
        vórtice central). Tratar $b_{fl}$ y $c_{fl}$ como dos parámetros globales
        efectivos evita importar números ajenos a nuestra geometría.
\end{itemize}
\end{antidoto}

\begin{lapizypapel}
\begin{enumerate}[leftmargin=1.5em]
  \item Resolvé paso a paso la ecuación de desaceleración libre
        $J\dot\omega + b\omega = -T_C$ y explicá cómo distinguir gráficamente si
        domina la fricción seca o la viscosa.
  \item Justificá físicamente por qué la carga del fluido tiene un término
        proporcional a $\omega_2$ y otro proporcional a $\omega_2|\omega_2|$
        (y por qué lleva valor absoluto).
  \item Explicá, mirando las ecuaciones de $J_1$ y $J_2$, por qué un sensor
        que sólo mide $\omega_1$ no detecta que la cápsula se desacopló.
\end{enumerate}
\end{lapizypapel}

% =====================================================================
\chapter{Medición sin sensores, zona ciega, arranque ciego y retardo}\label{cap:sensorless}
% =====================================================================

\begin{ideafisica}
En un motor BLDC \emph{sensorless} no hay ningún tacómetro ni encoder en el eje.
La velocidad se estima «escuchando» cuándo la tensión inducida (BEMF) de la fase
desconectada cruza por cero. De ese principio físico nacen dos limitaciones
enormes que condicionan todo el control: \textbf{por debajo de cierta velocidad
el sensor queda completamente ciego} (lo que obliga a arrancar a lazo abierto),
y \textbf{por encima de ella la medición siempre llega atrasada} en proporción
inversa a la velocidad.
\end{ideafisica}

En los capítulos anteriores construimos la planta física (motor + dos inercias +
acople + fluido). Ahora debemos modelar con la misma honestidad \emph{cómo
medimos la velocidad} y \emph{cómo logramos poner el motor en marcha desde el
reposo}.

% ---------------------------------------------------------------------
\section{Cómo se mide velocidad sin encoder: los cruces por cero}
% ---------------------------------------------------------------------

Recordemos la Figura~\ref{fig:bldc_estrella} del Capítulo~\ref{cap:motor}: en
conmutación de seis pasos, en todo momento hay dos fases llevando corriente y
una tercera fase abierta ($i_k = 0$). ¿Qué tensión aparece en el borne de esa
fase abierta respecto del neutro $n$? Como $i_k = 0$ y $\dot i_k = 0$, tanto la
caída resistiva $R\,i_k$ como la inductiva $L_s\dot i_k$ valen cero, y queda
limpia únicamente la BEMF:
\[
  v_{kn}(t) = e_k(t).
\]

Durante esa ventana de \SI{60}{\degree} eléctricos en que la fase queda libre,
su BEMF pasa de positiva a negativa (o de negativa a positiva), atravesando el
cero justo en el punto medio de la ventana. En cada ciclo eléctrico completo
($2\pi$ radianes eléctricos), una misma fase cruza por cero \textbf{dos veces}
(una subiendo y una bajando), separadas exactamente medio ciclo eléctrico
($\Delta\theta_e = \pi$).

Como $\theta_e = p\,\theta_1$, el tiempo $T_{ZC}$ entre dos cruces por cero
consecutivos de una misma fase es:
\begin{equation}
  p\,\omega_1\,T_{ZC} = \pi
  \quad\Longrightarrow\quad
  T_{ZC} = \frac{\pi}{p\,\omega_1} = \frac{30}{p\,n}\qquad [\si{s}],
  \label{eq:tzc}
\end{equation}
donde usamos $\omega_1 = \frac{\pi}{30}n$ para pasar de \si{rad/s} a \si{rpm}.
Un circuito comparador detecta cada cruce y un cronómetro mide el período
$T_{ZC}$; invirtiendo~\eqref{eq:tzc} se calcula la velocidad estimada:
\begin{equation}
  \hat\omega_1 = \frac{\pi}{p\,T_{ZC}},
  \qquad\text{o bien}\qquad
  \hat n = \frac{30}{p\,T_{ZC}}.
  \label{eq:omega_de_tzc}
\end{equation}

% ---------------------------------------------------------------------
\section{Segunda no linealidad dominante: la zona ciega a baja velocidad}
% ---------------------------------------------------------------------

Volvamos a la ley de Faraday del Capítulo~\ref{cap:motor}: la amplitud de la
tensión inducida es \textbf{proporcional a la velocidad}:
\[
  |e_{\text{pico}}| = k_e\,\omega_1.
\]
¿Qué sucede cuando el motor gira muy despacio ($\omega_1 \to 0$)? La señal de
BEMF se achica hasta volverse minúscula, mientras que el ruido eléctrico de
conmutación PWM, las histéresis y el umbral efectivo $V_{\text{umbral}}$ del
comparador siguen estando ahí.

Por debajo de una velocidad crítica
\begin{equation}
  \omega_{\text{ciego}} \approx \frac{V_{\text{umbral}}}{k_e}
  \qquad (n_{\text{ciego}} \approx 300\text{--}\SI{400}{rpm}\text{ en motores pequeños}),
  \label{eq:zona_ciega}
\end{equation}
la BEMF queda enterrada bajo el ruido (Figura~\ref{fig:bemf_ceguera}). El
comparador deja de entregar cruces válidos: \textbf{el error de medición salta
al \SI{100}{\percent}}. No estamos ante un pequeño ruido aditivo gaussiano, sino
ante la \textbf{desaparición total de la señal de realimentación}. Si pasa un
tiempo límite $\tau_{\text{stall}}$ sin recibir ningún cruce válido, el sistema
debe declarar pérdida de sincronismo (\emph{stall}) y reiniciar.

\begin{figure}[htbp]
\centering
\begin{minipage}[t]{0.48\linewidth}
\centering
\begin{tikzpicture}
\begin{axis}[
  width=6.6cm, height=4.6cm,
  axis lines=middle,
  xlabel={$t$}, ylabel={$e_k(t)$},
  xmin=0, xmax=6.5, ymin=-1.3, ymax=1.3,
  xtick=\empty, ytick=\empty,
  label style={font=\footnotesize},
  title={\footnotesize\bfseries Velocidad alta ($n > n_{\text{ciego}}$)},
  title style={yshift=-4pt}
]
  % Banda de ruido/umbral
  \fill[crimson!12] (axis cs:0,-0.22) rectangle (axis cs:6.5,0.22);
  \draw[crimson, dashed, thin] (axis cs:0,0.22) -- (axis cs:6.5,0.22);
  \draw[crimson, dashed, thin] (axis cs:0,-0.22) -- (axis cs:6.5,-0.22);
  \node[crimson, font=\tiny, above] at (axis cs:5.2,0.28) {$\pm V_{\text{umbral}}$};

  % Señal BEMF fuerte
  \addplot[inkblue, very thick, domain=0:6.28, samples=120] {1.05*sin(deg(x))};
  \fill[tealaccent] (axis cs:3.1416,0) circle (2.6pt);
  \draw[->, tealaccent, thin] (axis cs:2.25,-0.70) -- (axis cs:3.04,-0.08);
  \node[tealaccent, font=\scriptsize\bfseries, below] at (axis cs:1.85,-0.68) {Cruce claro};
\end{axis}
\end{tikzpicture}
\end{minipage}
\hfill
\begin{minipage}[t]{0.48\linewidth}
\centering
\begin{tikzpicture}
\begin{axis}[
  width=6.6cm, height=4.6cm,
  axis lines=middle,
  xlabel={$t$}, ylabel={$e_k(t)$},
  xmin=0, xmax=6.5, ymin=-1.3, ymax=1.3,
  xtick=\empty, ytick=\empty,
  label style={font=\footnotesize},
  title={\footnotesize\bfseries Velocidad baja ($n < n_{\text{ciego}}$)},
  title style={yshift=-4pt}
]
  % Banda de ruido/umbral
  \fill[crimson!15] (axis cs:0,-0.22) rectangle (axis cs:6.5,0.22);
  \draw[crimson, dashed, thin] (axis cs:0,0.22) -- (axis cs:6.5,0.22);
  \draw[crimson, dashed, thin] (axis cs:0,-0.22) -- (axis cs:6.5,-0.22);

  % Señal BEMF diminuta ahogada en ruido
  \addplot[mutedgray, thick, domain=0:6.28, samples=120] {0.14*sin(deg(0.5*x))};
  \node[crimson, font=\scriptsize\bfseries, align=center] at (axis cs:3.25,0.75) {BEMF tapada por\\umbral/ruido $\Rightarrow$ ceguera};
\end{axis}
\end{tikzpicture}
\end{minipage}
\caption{Origen físico de la zona ciega: como la amplitud de la BEMF escala con la velocidad, a bajas revoluciones queda sumergida dentro del umbral y ruido del comparador.}
\label{fig:bemf_ceguera}
\end{figure}

% ---------------------------------------------------------------------
\section{Física del arranque ciego a lazo abierto}
% ---------------------------------------------------------------------

La zona ciega plantea una paradoja inmediata: si en reposo ($\omega_1 = 0$) el
sensor no ve nada, \textbf{¿cómo hacemos para que el motor empiece a girar y
supere $n_{\text{ciego}}$?}

La respuesta es operar el motor inicialmente como un \textbf{motor síncrono a
lazo abierto}: sin mirar la posición del rotor, aplicamos a las tres bobinas del
estator tensiones senoidales desfasadas \SI{120}{\degree} entre sí, cuya
amplitud $V_{ph}(t)$ y frecuencia eléctrica $f_e(t)$ suben suavemente en rampa:
\begin{equation}
  v_a = V_{ph}(t)\cos(\theta_s), \quad
  v_b = V_{ph}(t)\cos\!\left(\theta_s - \tfrac{2\pi}{3}\right), \quad
  v_c = V_{ph}(t)\cos\!\left(\theta_s + \tfrac{2\pi}{3}\right),
  \label{eq:tensiones_arranque}
\end{equation}
donde $\theta_s(t) = 2\pi\int_0^t f_e(\tau)\,\dd\tau$ es el ángulo del campo
magnético rotante que estamos creando artificialmente en el estator.

\begin{ideafisica}[La asombrosa unidad física: dos acoples magnéticos en serie]
Mirá la Figura~\ref{fig:campo_rotante}. Durante el arranque ciego, el campo
magnético giratorio del estator (en el ángulo $\theta_s$) arrastra a los imanes
del rotor (en el ángulo eléctrico $\theta_e = p\theta_1$) mediante el
\textbf{exactamente mismo principio físico} con el que el imán inferior arrastra
a la cápsula agitadora! En ambos casos hay un ángulo de desfase y un torque
restaurador proporcional al seno de ese desfase.
\end{ideafisica}

\begin{figure}[htbp]
\centering
\begin{tikzpicture}[
  >={Stealth[length=5pt]},
  thick
]
  % Círculo del estator
  \draw[inkblue!35, very thick] (0,0) circle (2.0cm);
  \draw[mutedgray, dashed] (-2.3,0) -- (2.3,0);
  \draw[mutedgray, dashed] (0,-2.3) -- (0,2.3);

  % Vector rotor (ángulo theta_e = 25 deg)
  \draw[->, tealaccent, ultra thick] (0,0) -- (22:1.65cm)
    node[below right, font=\small\bfseries]{Rotor ($\theta_e = p\theta_1$)};

  % Vector campo estator (ángulo theta_s = 72 deg)
  \draw[->, inkblue, ultra thick] (0,0) -- (72:1.95cm)
    node[above right, font=\small\bfseries]{Campo estator ($\theta_s$)};

  % Arco de ángulo de carga gamma = theta_s - theta_e
  \draw[->, amberwarn, very thick] (22:1.05cm) arc (22:72:1.05cm);
  \node[amberwarn, font=\small\bfseries] at (47:1.35cm) {$\gamma$};

  % Leyenda al costado
  \node[right, align=left, font=\footnotesize] at (4.4,0.35) {
    Ángulo de carga eléctrico:\\[2pt]
    $\gamma = \theta_s(t) - p\,\theta_1(t)$\\[6pt]
    $\bullet\; 0 < \gamma < 90^\circ$: torque arrastra al rotor.\\[2pt]
    $\bullet\; \gamma > 90^\circ$: pierde sincronismo!
  };
\end{tikzpicture}
\caption{Durante el arranque ciego, el campo giratorio del estator ($\theta_s$) tira del imán del rotor ($\theta_e = p\theta_1$) con un torque proporcional a $\sin(\theta_s - p\theta_1)$.}
\label{fig:campo_rotante}
\end{figure}

\subsection{Deducción limpia del torque de arranque (en 6 líneas)}
Para que la fórmula del arranque no aparezca nunca «de la galera», deduzcámosla
paso a paso. A las frecuencias muy bajas del arranque ($f_e < \SI{10}{Hz}$), la
reactancia inductiva es despreciable ($2\pi f_e L_s \ll R$), de modo que en cada
fase $k \in \{0,1,2\}$ (con ángulos $\phi_k = \frac{2\pi k}{3}$) vale simplemente
la ley de Ohm con BEMF senoidal:
\[
  e_k = -k_e\,\omega_1\,\sin(\theta_e - \phi_k),
  \qquad
  i_k = \frac{v_k - e_k}{R}
      = \frac{V_{ph}\cos(\theta_s - \phi_k) + k_e\,\omega_1\sin(\theta_e - \phi_k)}{R}.
\]
(El signo menos en $e_k$ proviene de la ley de Faraday $e_k = \frac{\dd}{\dd t}[\psi_m\cos(\theta_e - \phi_k)] = -p\psi_m\omega_1\sin(\theta_e-\phi_k)$).
Por conservación de energía, la potencia electromecánica total de las tres fases
es $P_{\text{em}} = \sum_{k=0}^2 e_k i_k = T_{e,\text{arr}}\,\omega_1$.
Dividiendo por $\omega_1$:
\[
  T_{e,\text{arr}}
  = -\frac{k_e}{R}\sum_{k=0}^{2}
    \Big[V_{ph}\sin(\theta_e - \phi_k)\cos(\theta_s - \phi_k)
       + k_e\omega_1\sin^2(\theta_e - \phi_k)\Big].
\]
Usamos dos identidades trigonométricas elementales y el hecho de que tres
vectores iguales separados \SI{120}{\degree} suman cero:
\begin{align*}
  \sum_{k=0}^2 \sin(\theta_e-\phi_k)\cos(\theta_s-\phi_k)
  &= \frac{1}{2}\sum_{k=0}^2\Big[\sin(\theta_e-\theta_s) + \underbrace{\sin(\theta_e+\theta_s-2\phi_k)}_{\text{suma}=0}\Big]
   = -\frac{3}{2}\sin(\theta_s - \theta_e), \\
  \sum_{k=0}^2 \sin^2(\theta_e-\phi_k)
  &= \sum_{k=0}^2 \frac{1 - \overbrace{\cos(2\theta_e - 2\phi_k)}^{\text{suma}=0}}{2}
   = \frac{3}{2}.
\end{align*}
Reemplazando ambas sumas obtenemos sin ningún misterio:
\begin{ecuacionclave}
\begin{equation}
  T_{e,\text{arr}}(t) \approx \frac{3\,k_e}{2R}
  \Big(V_{ph}(t)\,\sin\!\big(\theta_s(t) - p\,\theta_1\big) - k_e\,\omega_1\Big).
  \label{eq:torque_arranque_deducido}
\end{equation}
\end{ecuacionclave}

\begin{lectura}
Mirá qué clara resulta~\eqref{eq:torque_arranque_deducido} ahora que sabés de
dónde salió cada letra:
\begin{itemize}
  \item El factor $\frac{3}{2}$ viene simplemente de sumar tres fases senoidales
        desfasadas \SI{120}{\degree}.
  \item El primer término, $\frac{3k_e V_{ph}}{2R}\sin(\theta_s - p\theta_1)$,
        es un \textbf{torque de enganche síncrono}: mientras $|\theta_s - p\theta_1| < \pi/2$,
        arrastra al rotor para que siga al campo rotante a velocidad
        $\omega_1 \approx 2\pi f_e(t)/p$.
  \item El segundo término, $-\frac{3k_e^2}{2R}\omega_1$, vuelve a ser nuestro
        viejo conocido \textbf{amortiguamiento eléctrico por BEMF}.
  \item \textbf{Por qué la rampa debe subir a la vez frecuencia y amplitud ($V/f$):}
        el torque máximo disponible (en $\sin = 1$) vale
        $\frac{3k_e}{2R}(V_{ph} - k_e\omega_1)$. A medida que el rotor acelera,
        la BEMF $k_e\omega_1$ crece y le «come» tensión a $V_{ph}$; si no
        subiéramos $V_{ph}(t)$ junto con la frecuencia $f_e(t)$, el torque
        disponible se desplomaría antes de salir de la zona ciega.
  \item \textbf{Cuándo falla el arranque:} si la rampa acelera demasiado rápido
        o la carga mecánica exige más que el torque máximo disponible, el ángulo
        $\theta_s - p\theta_1$ pasa de \SI{90}{\degree} y el rotor pierde el paso.
\end{itemize}
\end{lectura}

Una vez que la rampa (eventualmente ayudada por un escalón breve de ciclo de
trabajo al pasar a seis pasos) logra empujar al rotor por encima de
$n_{\text{ciego}}$, el comparador comienza a entregar cruces por cero limpios y
consistentes: recién ahí se realiza la \textbf{transición (\emph{handoff}) al
lazo cerrado} y el controlador PI toma el mando.

% ---------------------------------------------------------------------
\section{Tercera no linealidad dominante: retardo variable con la velocidad (\texorpdfstring{$\tau \propto 1/n$}{tau proporcional a 1/n})}
% ---------------------------------------------------------------------

Supongamos que ya superamos la zona ciega ($n > n_{\text{ciego}}$) y el sensor
funciona bien. ¿Significa eso que conocemos la velocidad $\omega_1(t)$ en tiempo
real? \textbf{No.} Medir velocidad cronometrando el período entre eventos
introducirá siempre un \textbf{retardo de tiempo} que empeora a medida que baja
la velocidad.

Observemos la línea de tiempo de la Figura~\ref{fig:retardo_zc}. Existen tres
fuentes de atraso acumulativas:
\begin{enumerate}
  \item \textbf{Ventana de promedio móvil ($M$ períodos):} para suavizar pequeñas
        asimetrías entre fases o ruido, suele promediarse el tiempo de los
        últimos $M$ semiciclos eléctricos (duración total $M\,T_{ZC}$). La
        velocidad promediada corresponde en realidad al \emph{punto medio} de
        esa ventana pasada, es decir, ya nace atrasada $\frac{M}{2}T_{ZC}$.
  \item \textbf{Retención entre cruces (ZOH por eventos):} entre un cruce por
        cero y el siguiente no entra información nueva; el controlador sigue
        leyendo un valor congelado en escalera. En promedio, ese congelamiento
        agrega medio período más: $\frac{1}{2}T_{ZC}$.
  \item \textbf{Muestreo del controlador discreto:} si el controlador ejecuta
        cada $\Delta t_c$ segundos, su retención de orden cero aporta en promedio
        $\frac{1}{2}\Delta t_c$.
\end{enumerate}

\begin{figure}[htbp]
\centering
\begin{tikzpicture}[>={Stealth[length=5pt]}, thick]
  % Ejes
  \draw[->] (0,0) -- (11.2,0) node[right, font=\footnotesize]{$t$};
  \draw[->] (0,0) -- (0,3.3) node[above, font=\footnotesize]{Velocidad};

  % Velocidad real creciendo
  \draw[inkblue, very thick] (0.2,0.7) .. controls (3.0,1.0) and (6.5,2.1) .. (10.5,2.75);
  \node[inkblue, font=\footnotesize\bfseries] at (9.6,3.05) {$\omega_1(t)$ real};

  % Cruces por cero (más espaciados al principio, más juntos después)
  \foreach \x in {0.8, 3.0, 4.9, 6.5, 7.8, 8.9, 9.9} {
    \draw[mutedgray, dotted] (\x,0) -- (\x,2.85);
    \fill[amberwarn] (\x,0) circle (2.2pt);
  }
  \node[below, amberwarn, font=\scriptsize] at (1.9,-0.08) {$T_{ZC}$ largo (baja rpm)};
  \node[below, amberwarn, font=\scriptsize] at (8.35,-0.08) {$T_{ZC}$ corto (alta rpm)};

  % Señal medida escalonada y atrasada
  \draw[crimson, very thick]
    (0.8,0.72) -- (3.0,0.72) -- (3.0,0.85) -- (4.9,0.85) -- (4.9,1.18)
    -- (6.5,1.18) -- (6.5,1.58) -- (7.8,1.58) -- (7.8,1.95)
    -- (8.9,1.95) -- (8.9,2.25) -- (9.9,2.25) -- (9.9,2.48) -- (10.6,2.48);
  \node[crimson, font=\footnotesize\bfseries] at (9.4,1.65) {$\hat\omega_1(t)$ medida};

  % Flecha de retardo tau
  \draw[<->, amberwarn, very thick] (4.9,1.38) -- (6.05,1.38);
  \node[fill=white, inner sep=1.2pt, text=amberwarn, font=\scriptsize\bfseries] at (5.48,1.08) {$\tau(n)$};
\end{tikzpicture}
\caption{Efecto del muestreo por eventos de cruce por cero: la velocidad estimada $\hat\omega_1$ forma una escalera atrasada respecto de $\omega_1(t)$. Al bajar la velocidad, $T_{ZC}$ se alarga y el retardo $\tau$ crece como $1/n$.}
\label{fig:retardo_zc}
\end{figure}

Sumando las tres contribuciones y sustituyendo $T_{ZC} = \frac{\pi}{p\,\omega_1} = \frac{30}{p\,n}$:
\begin{ecuacionclave}
\begin{equation}
  \tau(n) \approx \frac{M + 1}{2}\,T_{ZC} + \frac{\Delta t_c}{2}
          = \frac{(M + 1)\,\pi}{2\,p\,\omega_1} + \frac{\Delta t_c}{2}
          = \frac{15\,(M + 1)}{p\,n} + \frac{\Delta t_c}{2}.
  \label{eq:retardo_cap5}
\end{equation}
\end{ecuacionclave}

\paragraph{Por qué un retardo desestabiliza un lazo de control.}
En el dominio de la frecuencia, un retardo puro de tiempo $\tau$ tiene función
de transferencia $e^{-s\tau}$. Evaluado en $s = j\omega$:
\[
  \big|e^{-j\omega\tau}\big| = 1,
  \qquad
  \arg\!\big(e^{-j\omega\tau}\big) = -\omega\,\tau \quad [\si{rad}].
\]
Un retardo \textbf{no atenúa la ganancia} en ninguna frecuencia, pero \textbf{roba
margen de fase} proporcionalmente a $\omega\,\tau$. Si tu lazo cerrado, diseñado
a una velocidad nominal alta, tiene una frecuencia de cruce $\omega_c$ y dispone
de un margen de fase $\Delta\phi$ (en radianes) antes de volverse oscilatorio,
al bajar la consigna de velocidad $n$ el retardo $\tau(n)$ se agranda como
$1/n$. El lazo pierde estabilidad en cuanto $\omega_c\,\tau(n) \approx \Delta\phi$,
es decir, aproximadamente por debajo de:
\begin{equation}
  n_{\min,\text{estab}} \approx \frac{15\,(M + 1)\,\omega_c}{p\,\Delta\phi}.
  \label{eq:n_min_estab}
\end{equation}
Esto explica un hecho muy importante en la práctica: \textbf{un controlador PI
con ganancias fijas que funciona perfecto a \SI{2000}{rpm} puede empezar a
oscilar al bajar la consigna a \SI{500}{rpm}}, incluso antes de tocar la zona
ciega.

\begin{antidoto}
En este capítulo capturamos las dos propiedades que realmente dominan el sensor
(ceguera debajo de $n_{\text{ciego}}$ y retardo $\tau \propto 1/n$). Dejamos
afuera efectos de segundo orden como la cuantización entera del contador de
tiempo (error $\propto n^2$, despreciable a velocidades de agitador), el
\emph{jitter} fino por ruido térmico, el retardo de fase de los filtros RC
analógicos a muy alta velocidad y los tiempos de enmascaramiento
(\emph{blanking}).
\end{antidoto}

\begin{lapizypapel}
\begin{enumerate}[leftmargin=1.5em]
  \item Explicá físicamente por qué en un BLDC \emph{sensorless} existe una
        velocidad mínima $n_{\text{ciego}}$ por debajo de la cual el error de
        medición es del \SI{100}{\percent}.
  \item En la ecuación del arranque ciego~\eqref{eq:torque_arranque_deducido},
        ¿qué significa el ángulo $(\theta_s - p\theta_1)$ y por qué la rampa
        tiene que subir simultáneamente frecuencia y amplitud de tensión?
  \item Deducí la fórmula del retardo $\tau(n) \approx \frac{15(M+1)}{p\,n} + \frac{\Delta t_c}{2}$
        sumando sus tres contribuciones, y explicá con margen de fase por qué un
        PI estable a alta velocidad puede oscilar a baja velocidad.
\end{enumerate}
\end{lapizypapel}

% =====================================================================
\chapter{Ensamble del modelo no lineal y punto de equilibrio}\label{cap:ensamble}
% =====================================================================

\begin{ideafisica}
Llegó el momento de juntar en una sola hoja las cuatro piezas que construimos
por separado. No vamos a introducir ninguna ley nueva: sólo vamos a escribir
juntas las ecuaciones de Kirchhoff, Newton, acople y cargas, y aprender a
calcular el \textbf{punto de equilibrio} siguiendo la cadena física desde el
fluido hacia la fuente.
\end{ideafisica}

% ---------------------------------------------------------------------
\section{El sistema completo de ecuaciones diferenciales}
% ---------------------------------------------------------------------

Una vez superado el arranque ciego, en operación normal a lazo cerrado el estado
físico continuo del sistema queda completamente descripto por \textbf{cuatro
variables de estado}:
\[
  x = \begin{bmatrix} I & \omega_1 & \omega_2 & \delta \end{bmatrix}^{\mathsf{T}},
\]
que son respectivamente la corriente eléctrica $I$, la velocidad del motor
$\omega_1$, la velocidad de la cápsula agitadora $\omega_2$ y el desfase angular
$\delta = \theta_1 - \theta_2$ en el acople magnético. La entrada de control es
el ciclo de trabajo saturado $d \in [d_{\min}, 1]$ (que impone la tensión media
$V = d\,V_{dc}$), y la única variable accesible al sensor \emph{sensorless} es
$\omega_1$ (a través de cruces por cero, con zona ciega y retardo).

Reuniendo~\eqref{eq:bldc_sixstep}, \eqref{eq:newton1}--\eqref{eq:cinematica_delta},
\eqref{eq:acople_seno}, \eqref{eq:friccion_cap4} y~\eqref{eq:carga_fluido}:
\begin{ecuacionclave}
\begin{align}
  2L_s\,\frac{\dd I}{\dd t}        &= d\,V_{dc} - 2R\,I - 2k_e\,\omega_1,
  &&\text{(malla eléctrica)} \label{eq:nl_1}\\[3pt]
  J_1\,\frac{\dd\omega_1}{\dd t}   &= 2k_e\,I - T_{\max}\sin(N_a\delta)
                                    - T_{C1}\tanh\!\left(\frac{\omega_1}{\omega_{0,f}}\right) - b_1\omega_1,
  &&\text{(rotor)} \label{eq:nl_2}\\[3pt]
  J_2\,\frac{\dd\omega_2}{\dd t}   &= T_{\max}\sin(N_a\delta) - b_{fl}\omega_2 - c_{fl}\omega_2|\omega_2|
                                    - T_{C2}\tanh\!\left(\frac{\omega_2}{\omega_{0,f}}\right) - b_2\omega_2,
  &&\text{(agitador)} \label{eq:nl_3}\\[3pt]
  \frac{\dd\delta}{\dd t}          &= \omega_1 - \omega_2.
  &&\text{(cinemática)} \label{eq:nl_4}
\end{align}
\end{ecuacionclave}
Si hace falta generar los cruces por cero del sensor virtual o el ángulo del
arranque ciego, se agrega la relación trivial $\dot\theta_1 = \omega_1$.

Observá qué compacto es todo el sistema: \textbf{cuatro ecuaciones diferenciales
acopladas}, cada una con significado físico transparente. El cuadro~\ref{tab:resumen_nl}
resume dónde aparecen las cuatro no linealidades dominantes del problema.

\begin{table}[htbp]
\centering
\small
\begin{tabular}{@{}cl>{\raggedright\arraybackslash}p{6.8cm}@{}}
\toprule
\# & No linealidad dominante & Dónde vive y qué hace físicamente \\
\midrule
1 & Saturación de tensión & $d \in [d_{\min}, 1]$ en~\eqref{eq:nl_1}: limita corriente y aceleración máximas. \\
2 & Acople magnético no lineal & $T_{\max}\sin(N_a\delta)$ en~\eqref{eq:nl_2}--\eqref{eq:nl_3}: ablandamiento y desenganche al superar $T_{\max}$ o $\delta_{\text{crit}}$. \\
3 & Arrastre cuadrático y fricción seca & $c_{fl}\omega_2|\omega_2|$ y $T_{C,i}\tanh(\cdot)$ en~\eqref{eq:nl_2}--\eqref{eq:nl_3}: la carga crece cuadráticamente con la velocidad. \\
4 & Sensor BEMF (ceguera + retardo) & Ceguera para $n_1 < n_{\text{ciego}}$ (obliga a arranque ciego) y retardo $\tau \propto 1/n_1$ en la realimentación. \\
\bottomrule
\end{tabular}
\caption{Resumen de las no linealidades dominantes incluidas en el modelo.}
\label{tab:resumen_nl}
\end{table}

% ---------------------------------------------------------------------
\section{Cálculo del punto de equilibrio (de la carga hacia la fuente)}
% ---------------------------------------------------------------------

Antes de linealizar cualquier sistema no lineal hay que encontrar un
\textbf{punto de equilibrio} (o punto de operación): un estado estacionario en
el que el sistema gira acoplado a una velocidad deseada constante $\omega_0 > 0$
(por encima de la zona ciega y con $\omega_0 \gg \omega_{0,f}$, de modo que
$\tanh(\omega_0/\omega_{0,f}) \approx 1$).

Por definición, en equilibrio \textbf{todas las derivadas temporales valen cero}:
$\dot I = 0$, $\dot\omega_1 = 0$, $\dot\omega_2 = 0$ y $\dot\delta = 0$. El
secreto para resolverlo a mano sin enredarse es avanzar \textbf{de atrás para
adelante}, desde el fluido hacia la fuente eléctrica:

\begin{enumerate}
  \item \textbf{De la cinemática~\eqref{eq:nl_4} ($\dot\delta = 0$):}
        el motor y la cápsula giran a la misma velocidad:
        \begin{equation}
          \omega_1 = \omega_2 = \omega_0.
          \label{eq:eq_w}
        \end{equation}
  \item \textbf{Del agitador~\eqref{eq:nl_3} ($\dot\omega_2 = 0$):}
        el acople debe compensar exactamente la carga total sobre la cápsula
        $T_{L,\text{tot}}(\omega_0) = b_{fl}\omega_0 + c_{fl}\omega_0^2 + T_{C2} + b_2\omega_0$:
        \begin{equation}
          T_{\max}\sin(N_a\delta^\ast) = T_{L,\text{tot}}(\omega_0)
          \quad\Longrightarrow\quad
          \delta^\ast = \frac{1}{N_a}\arcsin\!\left(\frac{T_{L,\text{tot}}(\omega_0)}{T_{\max}}\right).
          \label{eq:eq_delta}
        \end{equation}
        \emph{Primera condición de existencia:} sólo existe equilibrio acoplado
        si $T_{L,\text{tot}}(\omega_0) < T_{\max}$ (y elegimos la rama estable
        $0 \le \delta^\ast < \delta_{\text{crit}}$).
  \item \textbf{Del rotor~\eqref{eq:nl_2} ($\dot\omega_1 = 0$):}
        el motor debe proveer el torque que se lleva el acople ($T_{L,\text{tot}}$)
        más su propia fricción ($T_{f1}(\omega_0) = T_{C1} + b_1\omega_0$):
        \begin{equation}
          2k_e\,I_0 = T_{L,\text{tot}}(\omega_0) + T_{C1} + b_1\omega_0
          \quad\Longrightarrow\quad
          I_0 = \frac{T_{L,\text{tot}}(\omega_0) + T_{C1} + b_1\omega_0}{2k_e}.
          \label{eq:eq_I}
        \end{equation}
  \item \textbf{De la malla eléctrica~\eqref{eq:nl_1} ($\dot I = 0$):}
        la tensión aplicada debe cubrir la caída óhmica y la BEMF:
        \begin{equation}
          d_0\,V_{dc} = 2R\,I_0 + 2k_e\,\omega_0
          \quad\Longrightarrow\quad
          d_0 = \frac{2R\,I_0 + 2k_e\,\omega_0}{V_{dc}}.
          \label{eq:eq_d}
        \end{equation}
        \emph{Segunda condición de existencia:} sólo es alcanzable si
        $d_0 \in [d_{\min},\, 1]$.
\end{enumerate}

Fijate qué natural resulta: no hubo que resolver ningún sistema de ecuaciones
enredado. Cada ecuación se despeja en cascada a partir de la anterior, y en el
camino saltan a la vista las dos fronteras físicas estáticas del sistema
($T_{L,\text{tot}} < T_{\max}$ y $d_0 \le 1$).

% ---------------------------------------------------------------------
\section{Ejemplo numérico ilustrativo hecho a mano}\label{sec:ejemplo_numerico}
% ---------------------------------------------------------------------

Para fijar los órdenes de magnitud reales de un agitador pequeño con motor de
disco rígido (sin depender de ninguna medición definitiva del prototipo, sino
usando valores típicos físicamente consistentes), hagamos una cuenta completa a
mano. Supongamos:
\begin{itemize}
  \item \textbf{Motor y fuente:} $V_{dc} = \SI{12}{V}$, $p = 2$,
        $R = \SI{2,0}{\ohm}$ ($2R = \SI{4,0}{\ohm}$),
        $L_s = \SI{1,0}{mH}$ ($2L_s = \SI{2,0}{mH}$),
        $k_e = \SI{0,012}{V.s/rad}$ ($2k_e = \SI{0,024}{N.m/A}$).
  \item \textbf{Inercias y fricciones:}
        $J_1 = \SI{2,0e-5}{kg.m^2}$, $J_2 = \SI{2,0e-6}{kg.m^2}$;
        $T_{C1} = \SI{0,50}{mN.m}$, $b_1 = \SI{2,0e-6}{N.m.s/rad}$;
        $T_{C2} = \SI{0,30}{mN.m}$, $b_2 = \SI{1,0e-6}{N.m.s/rad}$.
  \item \textbf{Acople y fluido (agua):}
        $N_a = 1$ ($\delta_{\text{crit}} = \SI{90}{\degree}$),
        $T_{\max} = \SI{6,0}{mN.m}$,
        $b_{fl} = \SI{8,0e-6}{N.m.s/rad}$,
        $c_{fl} = \SI{1,5e-7}{N.m.s^2/rad^2}$.
\end{itemize}

Calculemos el punto de equilibrio para una consigna de
$n_0 = \SI{1000}{rpm}$, es decir:
\[
  \omega_0 = \frac{\pi}{30}\cdot 1000 \approx \SI{104,7}{rad/s}.
\]
\begin{enumerate}
  \item \textbf{Carga total sobre la cápsula:}
        \begin{align*}
          b_{fl}\,\omega_0 + b_2\,\omega_0 &= (8{,}0 + 1{,}0)\times 10^{-6}\times 104{,}7 = \SI{0,94}{mN.m},\\
          c_{fl}\,\omega_0^2 &= 1{,}5\times 10^{-7}\times (104{,}7)^2 = \SI{1,65}{mN.m},\\
          T_{L,\text{tot}}(\omega_0) &= 0{,}94 + 1{,}65 + 0{,}30 = \SI{2,89}{mN.m}.
        \end{align*}
        Como $\SI{2,89}{mN.m} < T_{\max} = \SI{6,0}{mN.m}$ (casi la mitad del
        máximo), la cápsula gira acoplada con buen margen.
  \item \textbf{Desfase del acople:}
        \[
          \delta^\ast = \arcsin\!\left(\frac{2{,}89}{6{,}0}\right)
                      = \arcsin(0{,}481)
                      \approx \SI{0,502}{rad} \approx \SI{28,8}{\degree}.
        \]
  \item \textbf{Corriente estacionaria:} sumamos la fricción del rotor
        $T_{f1}(\omega_0) = 0{,}50 + 2{,}0\times 10^{-6}\times 104{,}7 = \SI{0,71}{mN.m}$:
        \[
          T_{e,0} = 2{,}89 + 0{,}71 = \SI{3,60}{mN.m}
          \quad\Longrightarrow\quad
          I_0 = \frac{3{,}60\times 10^{-3}}{0{,}024} = \SI{0,150}{A} = \SI{150}{mA}.
        \]
  \item \textbf{Tensión y ciclo de trabajo:}
        \[
          V_0 = \underbrace{4{,}0\times 0{,}150}_{\SI{0,60}{V}\text{ (cobre)}}
              + \underbrace{0{,}024\times 104{,}7}_{\SI{2,51}{V}\text{ (BEMF)}}
              = \SI{3,11}{V}
          \quad\Longrightarrow\quad
          d_0 = \frac{3{,}11}{12{,}0} \approx 0{,}259 \;\; (\SI{25,9}{\percent}).
        \]
\end{enumerate}

Mirá toda la información física que sacamos en cuatro renglones: a \SI{1000}{rpm},
el \SI{81}{\percent} de la tensión se va en BEMF y sólo el \SI{19}{\percent} en
resistencia; el motor consume \SI{150}{mA}, tiene mucho margen de tensión
($d_0 \approx 0{,}26$), y la cápsula va atrasada unos \SI{29}{\degree} respecto
del imán inferior. En cambio, si subiéramos la velocidad hasta alrededor de
\SI{1600}{rpm} ($\omega \approx \SI{167}{rad/s}$), el término cuadrático
$c_{fl}\omega^2$ llevaría $T_{L,\text{tot}}$ hasta \SI{6,0}{mN.m} y la cápsula
se desacoplaría aunque el motor todavía tuviese tensión de sobra ($d < 0{,}45$).

\begin{lapizypapel}
\begin{enumerate}[leftmargin=1.5em]
  \item Escribí sin mirar el texto las cuatro ecuaciones diferenciales
        \eqref{eq:nl_1}--\eqref{eq:nl_4} y verificá el signo de cada término.
  \item Si fijás una velocidad de equilibrio $\omega_0$, explicá en qué orden se
        despejan $\delta^\ast$, $I_0$ y $d_0$, y qué dos desigualdades físicas
        deben cumplirse para que ese equilibrio exista.
\end{enumerate}
\end{lapizypapel}

% =====================================================================
\chapter{Linealización paso a paso: de las ecuaciones no lineales a \texorpdfstring{$G(s)$}{G(s)}}\label{cap:linealizacion}
% =====================================================================

\begin{ideafisica}
Las herramientas clásicas de control (polos, ceros, lugar de raíces, Bode,
diseño analítico de un PI) funcionan sobre sistemas \textbf{lineales}.
Linealizar no es «hacer trampa»: es poner una lupa sobre un punto de operación
$(I_0, \omega_0, \delta^\ast, d_0)$ y reemplazar cada curva suave por su
\textbf{recta tangente} en ese punto, sabiendo de antemano que la aproximación
sólo vale para desviaciones pequeñas alrededor de él.
\end{ideafisica}

En este capítulo vamos a derivar \textbf{término por término y coeficiente por
coeficiente} tanto la representación en espacio de estados $(A, B, C)$ como las
funciones de transferencia $G_2(s)$ (dos inercias) y $G_1(s)$ (acople rígido).
El objetivo es que, si alguien te señala cualquier casilla de la matriz $A$,
puedas decirle al instante de qué ley física proviene y por qué lleva ese signo.

% ---------------------------------------------------------------------
\section{Qué significa linealizar: variables incrementales y serie de Taylor}
% ---------------------------------------------------------------------

Si una función no lineal suave $f(x)$ se evalúa muy cerca de un punto de
equilibrio $x_0$, definimos la \textbf{variable incremental} (la desviación
pequeña respecto del equilibrio):
\[
  \tilde x(t) = x(t) - x_0 \qquad\Longleftrightarrow\qquad x(t) = x_0 + \tilde x(t).
\]
Como $x_0$ es constante, su derivada temporal es nula y resulta
$\dot{\tilde x} = \dot x$. Aplicando el desarrollo de Taylor de primer orden
(la ecuación de la recta tangente, Figura~\ref{fig:taylor}):
\begin{equation}
  f(x_0 + \tilde x) \approx f(x_0) + \underbrace{\left.\frac{\dd f}{\dd x}\right|_{x_0}}_{\text{pendiente local}}\tilde x.
  \label{eq:taylor_1}
\end{equation}

\begin{figure}[htbp]
\centering
\begin{tikzpicture}
\begin{axis}[
  width=9.2cm, height=4.8cm,
  axis lines=left,
  xlabel={Variable $x$},
  ylabel={$f(x)$},
  xmin=0, xmax=2.3, ymin=0, ymax=2.3,
  xtick={1.1}, xticklabels={$x_0$},
  ytick={0.85}, yticklabels={$f(x_0)$},
  tick label style={font=\footnotesize},
  label style={font=\small},
  clip=false
]
  % Entorno de validez (dibujado primero para quedar de fondo)
  \fill[tealaccent!11] (axis cs:0.85,0) rectangle (axis cs:1.35,2.25);
  \node[text=tealaccent, font=\scriptsize, align=center] at (axis cs:1.10,1.98) {Entorno pequeño\\$|\tilde x|$ chico};

  % Curva no lineal (ej. 0.70*x^2)
  \addplot[inkblue, very thick, domain=0:1.8, samples=80] {0.70*x^2};
  \node[text=inkblue, font=\footnotesize\bfseries, right] at (axis cs:1.82,2.20) {Curva real $f(x)$};

  % Recta tangente en x0=1.1: f(1.1)=0.847, f'(1.1)=1.54
  \addplot[tealaccent, very thick, dashed, domain=0.45:1.95] {0.847 + 1.54*(x - 1.1)};
  \node[text=tealaccent, font=\footnotesize\bfseries, right] at (axis cs:1.96,1.88) {Aprox. lineal};

  \fill[crimson] (axis cs:1.1,0.847) circle (2.8pt);
  \draw[mutedgray, dotted] (axis cs:1.1,0) -- (axis cs:1.1,0.847) -- (axis cs:0,0.847);
\end{axis}
\end{tikzpicture}
\caption{Interpretación geométrica de la linealización: cerca de $x_0$ la curva $f(x)$ se confunde con su recta tangente; lejos de $x_0$ el error crece.}
\label{fig:taylor}
\end{figure}

Ahora apliquemos~\eqref{eq:taylor_1} a los \textbf{únicos tres términos no
lineales} de nuestras ecuaciones~\eqref{eq:nl_1}--\eqref{eq:nl_4} alrededor de
$(I_0, \omega_0, \omega_0, \delta^\ast, d_0)$ (con $\omega_0 \gg \omega_{0,f}$):

\begin{enumerate}
  \item \textbf{Acople magnético $T_c(\delta) = T_{\max}\sin(N_a\delta)$:}
        su derivada respecto de $\delta$ es $N_a T_{\max}\cos(N_a\delta)$, por
        lo que:
        \[
          T_{\max}\sin\big(N_a(\delta^\ast + \tilde\delta)\big)
          \approx T_{\max}\sin(N_a\delta^\ast)
          + \underbrace{N_a\,T_{\max}\cos(N_a\delta^\ast)}_{K_{\text{eff}}(\delta^\ast)}\,\tilde\delta.
        \]
  \item \textbf{Arrastre cuadrático del fluido $c_{fl}\omega_2|\omega_2|$:}
        como estamos girando hacia adelante ($\omega_0 > 0$), cerca del
        equilibrio $\omega_2|\omega_2| = \omega_2^2$, cuya derivada es
        $2\omega_2$:
        \[
          c_{fl}(\omega_0 + \tilde\omega_2)^2
          \approx c_{fl}\omega_0^2 + \underbrace{2\,c_{fl}\,\omega_0}_{\text{amort.\ incremental}}\,\tilde\omega_2.
        \]
        \emph{Atención al factor~2:} la pendiente de la parábola $c_{fl}\omega^2$
        es $2c_{fl}\omega_0$ (el doble que el cociente $c_{fl}\omega_0$). Cuanto
        más rápido gira el agitador, más «duro» siente el fluido ante pequeños
        cambios de velocidad.
  \item \textbf{Fricción seca $T_{C,i}\tanh(\omega_i/\omega_{0,f})$:}
        para $\omega_0 \gg \omega_{0,f}$, $\tanh$ está en su meseta plana igual
        a $+1$ con derivada prácticamente nula:
        \[
          T_{C,i}\tanh\!\left(\frac{\omega_0 + \tilde\omega_i}{\omega_{0,f}}\right)
          \approx T_{C,i} + 0\cdot\tilde\omega_i.
        \]
        \emph{Consecuencia clave:} la fricción seca $T_{C,i}$ \textbf{desaparece
        de las ecuaciones linealizadas} (su derivada es cero), aunque sí alteró
        el punto de equilibrio $(I_0, d_0, \delta^\ast)$!
\end{enumerate}

% ---------------------------------------------------------------------
\section{Construcción de las matrices de estado \texorpdfstring{$(A, B, C)$}{(A, B, C)} fila por fila}
% ---------------------------------------------------------------------

Sustituimos $I = I_0 + \tilde I$, $\omega_1 = \omega_0 + \tilde\omega_1$,
$\omega_2 = \omega_0 + \tilde\omega_2$, $\delta = \delta^\ast + \tilde\delta$ y
$d = d_0 + \tilde d$ en las cuatro ecuaciones~\eqref{eq:nl_1}--\eqref{eq:nl_4}.
Veamos cómo ocurre la cancelación en la ecuación del agitador~\eqref{eq:nl_3}
(en las otras tres sucede exactamente lo mismo):
\begin{align*}
  J_2\,\dot{\tilde\omega}_2
  &\approx \Big[T_{\max}\sin(N_a\delta^\ast) + K_{\text{eff}}\tilde\delta\Big]
         - b_{fl}(\omega_0 + \tilde\omega_2)
         - \Big[c_{fl}\omega_0^2 + 2c_{fl}\omega_0\tilde\omega_2\Big]
         - T_{C2} - b_2(\omega_0 + \tilde\omega_2) \\[4pt]
  &= \underbrace{\Big(T_{\max}\sin(N_a\delta^\ast) - b_{fl}\omega_0 - c_{fl}\omega_0^2 - T_{C2} - b_2\omega_0\Big)}_{=\,0\text{ por equilibrio~\eqref{eq:eq_delta}!}}
   + K_{\text{eff}}\,\tilde\delta
   - \underbrace{(b_2 + b_{fl} + 2c_{fl}\omega_0)}_{B_{2,\text{eq}}(\omega_0)}\,\tilde\omega_2.
\end{align*}

¡Todos los términos constantes de equilibrio se cancelan exactamente! Definiendo
los coeficientes de amortiguamiento incremental equivalentes:
\begin{equation}
  B_{1,\text{eq}} = b_1,
  \qquad
  B_{2,\text{eq}}(\omega_0) = b_2 + b_{fl} + 2\,c_{fl}\,\omega_0,
  \qquad
  K_{\text{eff}}(\delta^\ast) = N_a\,T_{\max}\cos(N_a\delta^\ast),
  \label{eq:coefs_lin}
\end{equation}
las cuatro ecuaciones diferenciales lineales en variables incrementales quedan:
\begin{align}
  \dot{\tilde I}        &= -\frac{R}{L_s}\,\tilde I - \frac{k_e}{L_s}\,\tilde\omega_1 + \frac{V_{dc}}{2L_s}\,\tilde d, \label{eq:lin_1}\\
  \dot{\tilde\omega}_1  &= \frac{2k_e}{J_1}\,\tilde I - \frac{B_{1,\text{eq}}}{J_1}\,\tilde\omega_1 - \frac{K_{\text{eff}}}{J_1}\,\tilde\delta, \label{eq:lin_2}\\
  \dot{\tilde\omega}_2  &= -\frac{B_{2,\text{eq}}}{J_2}\,\tilde\omega_2 + \frac{K_{\text{eff}}}{J_2}\,\tilde\delta, \label{eq:lin_3}\\
  \dot{\tilde\delta}    &= \tilde\omega_1 - \tilde\omega_2. \label{eq:lin_4}
\end{align}

Escribiendo el vector de estado incremental
$\tilde x = [\tilde I,\ \tilde\omega_1,\ \tilde\omega_2,\ \tilde\delta]^{\mathsf{T}}$
en la forma estándar $\dot{\tilde x} = A\,\tilde x + B\,\tilde d$ y salida de
velocidad medida $\tilde\omega_1 = C\,\tilde x$ (en \si{rad/s}; si se desea en
\si{rpm}, el único elemento no nulo de $C$ pasa de $1$ a $30/\pi$):
\begin{ecuacionclave}
\begin{equation}
  A = \begin{bmatrix}
    -\dfrac{R}{L_s} & -\dfrac{k_e}{L_s} & 0 & 0 \\[8pt]
    \dfrac{2k_e}{J_1} & -\dfrac{B_{1,\text{eq}}}{J_1} & 0 & -\dfrac{K_{\text{eff}}}{J_1} \\[8pt]
    0 & 0 & -\dfrac{B_{2,\text{eq}}}{J_2} & \dfrac{K_{\text{eff}}}{J_2} \\[8pt]
    0 & 1 & -1 & 0
  \end{bmatrix},
  \qquad
  B = \begin{bmatrix}
    \dfrac{V_{dc}}{2L_s} \\[8pt] 0 \\[6pt] 0 \\[6pt] 0
  \end{bmatrix},
  \qquad
  C = \begin{bmatrix} 0 & 1 & 0 & 0 \end{bmatrix}.
  \label{eq:ss_cap7}
\end{equation}
\end{ecuacionclave}

\begin{lectura}[Mapa de defensa oral: por qué cada entrada de $A$ es exactamente esa]
Si te detienen en un examen y te señalan cualquier entrada de $A$, leé fila
como \emph{«¿qué hace cambiar a esta variable?»} y columna como \emph{«por culpa
de cuál otra?»}:
\begin{itemize}
  \item \textbf{Fila~1 ($\dot{\tilde I}$):} $A_{11} = -R/L_s$ (la resistencia frena
        la corriente; el factor 2 de $2R$ y $2L_s$ se cancela); $A_{12} = -k_e/L_s$
        (la BEMF $2k_e\omega_1$ dividida por $2L_s$ se opone a la corriente);
        $A_{13} = A_{14} = 0$ (la bobina no «ve» directamente a la cápsula).
  \item \textbf{Fila~2 ($\dot{\tilde\omega}_1$):} $A_{21} = +2k_e/J_1$ (más corriente
        acelera al rotor); $A_{22} = -B_{1,\text{eq}}/J_1$ (rozamiento frena al
        rotor); $A_{23} = 0$ (sin amortiguamiento interno $c\dot\delta$, $\omega_2$
        no aparece directo); $A_{24} = -K_{\text{eff}}/J_1$ (si el rotor se
        adelanta, $\tilde\delta > 0$, el acople lo \textbf{tira hacia atrás}).
  \item \textbf{Fila~3 ($\dot{\tilde\omega}_2$):} $A_{33} = -B_{2,\text{eq}}/J_2$
        (fluido y rozamiento frenan la cápsula); $A_{34} = +K_{\text{eff}}/J_2$
        (si $\tilde\delta > 0$, el acople \textbf{empuja hacia adelante} a la
        cápsula: signo opuesto a $A_{24}$ por acción y reacción!).
  \item \textbf{Fila~4 ($\dot{\tilde\delta}$):} $A_{42} = +1$ y $A_{43} = -1$
        porque $\delta = \theta_1 - \theta_2 \implies \dot\delta = \omega_1 - \omega_2$.
  \item \textbf{Matriz $B$:} la tensión $\tilde d\,V_{dc}$ sólo entra en la
        malla eléctrica, dividida por $2L_s$: $B_{11} = V_{dc}/(2L_s)$.
\end{itemize}
\end{lectura}

% ---------------------------------------------------------------------
\section{Reducción para análisis a mano: de 4.º orden a \texorpdfstring{$G_2(s)$ y $G_1(s)$}{G2(s) y G1(s)}}
% ---------------------------------------------------------------------

Una matriz $4\times 4$ es comodísima para la computadora, pero incómoda para
pensar con lápiz y papel. Recordemos del Capítulo~\ref{cap:motor} que la
constante de tiempo eléctrica $\tau_e = L_s/R$ es cientos de veces menor que las
constantes mecánicas. Veamos cuánto vale en nuestro ejemplo del
Capítulo~\ref{cap:ensamble}:
\[
  \tau_e = \frac{L_s}{R} = \frac{\SI{1,0}{mH}}{\SI{2,0}{\ohm}} = \SI{0,50}{ms}.
\]
Si despreciamos $2L_s\dot{\tilde I} \approx 0$ en~\eqref{eq:lin_1}, la corriente
responde instantáneamente:
\begin{equation}
  \tilde I \approx \frac{V_{dc}\,\tilde d - 2k_e\,\tilde\omega_1}{2R}
  \quad\Longrightarrow\quad
  \tilde T_e = 2k_e\,\tilde I = \frac{k_e V_{dc}}{R}\,\tilde d - \underbrace{\frac{2k_e^2}{R}}_{B_e}\,\tilde\omega_1.
  \label{eq:cuasiestatica}
\end{equation}
Al sustituir $\tilde T_e$ en~\eqref{eq:lin_2}, el amortiguamiento eléctrico
$B_e = 2k_e^2/R$ se suma directamente al rozamiento mecánico del rotor. Llamando
$B_1' = B_{1,\text{eq}} + \frac{2k_e^2}{R}$, el sistema se reduce a tres
ecuaciones puramente mecánicas:
\begin{align*}
  J_1\,\dot{\tilde\omega}_1 + B_1'\,\tilde\omega_1 &= \frac{k_e V_{dc}}{R}\,\tilde d - K_{\text{eff}}\,\tilde\delta, \\
  J_2\,\dot{\tilde\omega}_2 + B_{2,\text{eq}}\,\tilde\omega_2 &= K_{\text{eff}}\,\tilde\delta, \\
  \dot{\tilde\delta} &= \tilde\omega_1 - \tilde\omega_2.
\end{align*}

\subsection{Paso a Laplace y obtención de \texorpdfstring{$G_2(s)$}{G2(s)} sin saltar ningún renglón}
Tomamos transformada de Laplace con condiciones iniciales nulas (reemplazando
cada derivada $\dd/\dd t$ por multiplicar por $s$):
\begin{align}
  (J_1 s + B_1')\,\Omega_1(s) &= \frac{k_e V_{dc}}{R}\,D(s) - K_{\text{eff}}\,\Delta(s), \label{eq:lap_1}\\
  (J_2 s + B_{2,\text{eq}})\,\Omega_2(s) &= K_{\text{eff}}\,\Delta(s), \label{eq:lap_2}\\
  s\,\Delta(s) &= \Omega_1(s) - \Omega_2(s). \label{eq:lap_3}
\end{align}
Ahora eliminamos $\Omega_2(s)$ y $\Delta(s)$ en tres pasos algebraicos limpios:
\begin{enumerate}
  \item De~\eqref{eq:lap_3} despejamos $\Delta(s) = \frac{\Omega_1(s) - \Omega_2(s)}{s}$
        y lo sustituimos en~\eqref{eq:lap_2}:
        \[
          (J_2 s + B_{2,\text{eq}})\,\Omega_2(s) = \frac{K_{\text{eff}}}{s}\big(\Omega_1(s) - \Omega_2(s)\big).
        \]
  \item Multiplicamos por $s$ y agrupamos $\Omega_2(s)$ a la izquierda:
        \[
          (J_2 s^2 + B_{2,\text{eq}} s + K_{\text{eff}})\,\Omega_2(s) = K_{\text{eff}}\,\Omega_1(s)
          \;\Longrightarrow\;
          \Omega_2(s) = \frac{K_{\text{eff}}}{J_2 s^2 + B_{2,\text{eq}} s + K_{\text{eff}}}\,\Omega_1(s).
        \]
  \item Reemplazamos esa expresión de $\Omega_2(s)$ en $\Delta(s) = (\Omega_1 - \Omega_2)/s$:
        \[
          \Delta(s)
          = \frac{1}{s}\left(1 - \frac{K_{\text{eff}}}{J_2 s^2 + B_{2,\text{eq}} s + K_{\text{eff}}}\right)\Omega_1(s)
          = \frac{J_2 s + B_{2,\text{eq}}}{J_2 s^2 + B_{2,\text{eq}} s + K_{\text{eff}}}\,\Omega_1(s).
        \]
        (Fijate cómo el $K_{\text{eff}}$ del numerador se cancela y el factor $s$
        se simplifica limpiamente!).
  \item Finalmente sustituimos $\Delta(s)$ en~\eqref{eq:lap_1}, sacamos
        denominador común y despejamos $G_2(s) = \Omega_1(s)/D(s)$:
\end{enumerate}

\begin{ecuacionclave}
\begin{equation}
  G_2(s) = \frac{\Omega_1(s)}{D(s)}
         = \frac{\dfrac{k_e V_{dc}}{R}\,\big(J_2 s^2 + B_{2,\text{eq}} s + K_{\text{eff}}\big)}
                {(J_1 s + B_1')(J_2 s^2 + B_{2,\text{eq}} s + K_{\text{eff}}) + K_{\text{eff}}(J_2 s + B_{2,\text{eq}})}.
  \label{eq:g2s_cap7}
\end{equation}
\end{ecuacionclave}

% ---------------------------------------------------------------------
\section{Los dos casos límite reveladores de \texorpdfstring{$G_2(s)$}{G2(s)}}
% ---------------------------------------------------------------------

Una función de transferencia como~\eqref{eq:g2s_cap7} sólo se comprende de
verdad cuando probás sus dos casos extremos físicos:

\paragraph{Límite A: Acople muy rígido ($K_{\text{eff}} \to \infty$, lejos del desenganche).}
Dividamos numerador y denominador de~\eqref{eq:g2s_cap7} por $K_{\text{eff}}$ y
hagamos $K_{\text{eff}} \to \infty$:
\[
  G_2(s)
  = \frac{\dfrac{k_e V_{dc}}{R}\left(\dfrac{J_2 s^2 + B_{2,\text{eq}}s}{K_{\text{eff}}} + 1\right)}
         {(J_1 s + B_1')\left(\dfrac{J_2 s^2 + B_{2,\text{eq}}s}{K_{\text{eff}}} + 1\right) + (J_2 s + B_{2,\text{eq}})}
  \;\xrightarrow{K_{\text{eff}}\to\infty}\;
  \frac{\dfrac{k_e V_{dc}}{R}}{(J_1 + J_2)s + (B_1' + B_{2,\text{eq}})}.
\]
Llamando $J_{\text{tot}} = J_1 + J_2$ y $B_{\text{tot}} = B_{1,\text{eq}} + B_{2,\text{eq}}(\omega_0)$
(recordando que $B_1' = B_{1,\text{eq}} + \frac{2k_e^2}{R}$) y normalizando el
término independiente del denominador a~1, el sistema colapsa a una planta
clásica de \textbf{primer orden}:
\begin{ecuacionclave}
\begin{equation}
  G_1(s) = \frac{\Omega_1(s)}{D(s)} = \frac{K_{dc}}{\tau_{\text{mec}}\,s + 1},
  \qquad
  K_{dc} = \frac{V_{dc}/(2k_e)}{1 + \dfrac{R\,B_{\text{tot}}}{2k_e^2}},
  \qquad
  \tau_{\text{mec}} = \frac{J_{\text{tot}}}{B_{\text{tot}} + \dfrac{2k_e^2}{R}}.
  \label{eq:g1s_cap7}
\end{equation}
\end{ecuacionclave}
Mirá la ganancia estática $K_{dc}$: si no hubiera rozamientos ni fluido
($B_{\text{tot}} = 0$), vale exactamente $V_{dc}/(2k_e)$, que es la velocidad de
vacío a $d=1$ que hallamos en el Capítulo~\ref{cap:motor}. Y la constante de
tiempo mecánica $\tau_{\text{mec}}$ es la inercia total dividida por el
amortiguamiento total (mecánico + fluido + eléctrico $2k_e^2/R$).

\paragraph{Límite B: Acople perdido ($K_{\text{eff}} \to 0$, en $\delta^\ast \to \delta_{\text{crit}}$ o tras el desenganche).}
Pongamos ahora $K_{\text{eff}} = 0$ en~\eqref{eq:g2s_cap7}:
\[
  G_2(s)\Big|_{K_{\text{eff}} = 0}
  = \frac{\dfrac{k_e V_{dc}}{R}\,\cancel{(J_2 s^2 + B_{2,\text{eq}} s)}}
         {(J_1 s + B_1')\,\cancel{(J_2 s^2 + B_{2,\text{eq}} s)}}
  = \frac{k_e V_{dc}/R}{J_1 s + B_1'}.
\]
\textbf{¡Cancelación polo--cero exacta!} El polinomio del agitador
$(J_2 s^2 + B_{2,\text{eq}} s)$ se cancela entre numerador y denominador: en la
teoría lineal de sistemas, eso significa literalmente \textbf{pérdida de
observabilidad} de la dinámica de la cápsula desde la salida $\omega_1$. El
motor queda viendo únicamente su propia inercia $J_1$ y su propio rozamiento.

\paragraph{Continuación del ejemplo numérico en $n_0 = \SI{1000}{rpm}$.}
Con los números del Capítulo~\ref{cap:ensamble}:
\begin{itemize}
  \item $K_{\text{eff}} = \sqrt{6{,}0^2 - 2{,}89^2}\;\si{mN.m/rad} \approx \SI{5,26}{mN.m/rad}$
        (se ablandó un \SI{12}{\percent} respecto del vacío);
  \item frecuencia torsional:
        $\omega_{n,\text{tors}} = \sqrt{K_{\text{eff}}/J_{\text{eq}}}
        = \sqrt{5{,}26\times 10^{-3}/1{,}82\times 10^{-6}} \approx \SI{53,8}{rad/s}$ ($\approx \SI{8,6}{Hz}$);
  \item amortiguamiento incremental:
        $B_{1,\text{eq}} = \SI{2,0e-6}{N.m.s/rad}$,
        $B_{2,\text{eq}} = (1{,}0 + 8{,}0 + 2\cdot 0{,}15\cdot 104{,}7)\times 10^{-6} = \SI{4,04e-5}{N.m.s/rad}$,
        $B_{\text{tot}} = \SI{4,24e-5}{N.m.s/rad}$,
        frente al amortiguamiento eléctrico $B_e = \frac{2(0{,}012)^2}{2{,}0} = \SI{1,44e-4}{N.m.s/rad}$;
  \item planta de primer orden $G_1(s)$:
        \[
          \tau_{\text{mec}} = \frac{2{,}2\times 10^{-5}}{(4{,}24 + 14{,}4)\times 10^{-5}}
                            \approx \SI{0,118}{s} = \SI{118}{ms},
          \qquad
          K_{dc} = \frac{500}{1 + 4{,}24/14{,}4} \approx \SI{386}{rad/s}.
        \]
\end{itemize}
Confirmamos con números la separación de escalas: $\tau_{\text{mec}} = \SI{118}{ms}$
es \textbf{236 veces más lenta} que $\tau_e = \SI{0,50}{ms}$!

\begin{lapizypapel}
\begin{enumerate}[leftmargin=1.5em]
  \item ¿Por qué la fricción seca $T_{C,i}$ no aparece en la matriz $A$, mientras
        que el término cuadrático $c_{fl}\omega_2^2$ sí aparece con un factor~2
        ($2c_{fl}\omega_0$)?
  \item Explicá qué significa físicamente el elemento $A_{24} = -K_{\text{eff}}/J_1$
        y por qué $A_{34} = +K_{\text{eff}}/J_2$ tiene el signo contrario.
  \item Mostrá con dos renglones de álgebra cómo $G_2(s)$ se reduce a $G_1(s)$
        cuando $K_{\text{eff}} \to \infty$, y qué significa físicamente la
        cancelación polo--cero cuando $K_{\text{eff}} \to 0$.
\end{enumerate}
\end{lapizypapel}

% =====================================================================
\chapter{Control lineal PI y rango de validez del modelo}\label{cap:control_validez}
% =====================================================================

\begin{ideafisica}
El objetivo de la metodología de control no es que el modelo lineal «sirva para
todo», sino exactamente lo contrario: \textbf{diseñar un controlador sencillo
sobre el modelo linealizado, probarlo sobre la planta no lineal y mapear con
precisión en qué fronteras físicas el diseño lineal deja de funcionar}.
\end{ideafisica}

Con la planta reducida $G_1(s)$ y el retardo del sensor $\tau(n)$ ya sobre la
mesa, cerremos el círculo: diseñemos el controlador PI a mano, entendamos por
qué necesita \emph{anti-windup} al chocar contra la saturación y tracemos el
mapa completo de validez.

% ---------------------------------------------------------------------
\section{Por qué un controlador PI (y por qué no P solo ni PID)}
% ---------------------------------------------------------------------

Definamos el error de velocidad angular $e_\omega(t) = \omega^\ast - \omega_1(t)$
(donde $\omega^\ast$ es la consigna). Sobre la planta lineal de primer orden
$G_1(s) = \frac{K_{dc}}{\tau_{\text{mec}}s + 1}$, comparemos tres acciones:

\paragraph{1. ¿Por qué no basta un control proporcional puro ($d = k_p\,e_\omega$)?}
Con un controlador $C(s) = k_p$, la función de transferencia de lazo cerrado es
$\frac{k_p K_{dc}}{\tau_{\text{mec}}s + 1 + k_p K_{dc}}$. Ante un escalón de
consigna $\omega^\ast$, aplicando el teorema del valor final, el error en
régimen permanente nunca llega a cero:
\[
  e_{\omega,\text{ss}} = \frac{1}{1 + k_p K_{dc}}\,\omega^\ast \neq 0.
\]
Físicamente es obvio: para sostener el motor girando contra la fricción y el
fluido hace falta un ciclo de trabajo $d_0 > 0$; si sólo multiplicás el error por
$k_p$, necesitás que sobre algo de error permanentemente para generar ese $d_0$.

\paragraph{2. Qué aporta la acción integral (controlador PI).}
Al sumar un integrador que acumula el error en el tiempo:
\begin{equation}
  d(t) = k_p\,e_\omega(t) + k_i\int_0^t e_\omega(\tau)\,\dd\tau
  \quad\Longleftrightarrow\quad
  C_\omega(s) = k_p + \frac{k_i}{s} = \frac{k_p s + k_i}{s},
  \label{eq:pi_continuo}
\end{equation}
el integrador sigue subiendo (o bajando) $d$ mientras exista la más mínima
diferencia entre $\omega^\ast$ y $\omega_1$. Sólo deja de moverse cuando
$e_\omega = 0$: el error estacionario frente a consignas constantes \emph{y}
frente a cambios constantes de carga (viscosidad o fricción) es
\textbf{exactamente cero}.

\paragraph{3. ¿Por qué omitimos la acción derivativa ($k_d = 0$)?}
Recordemos el Capítulo~\ref{cap:sensorless}: el sensor de cruces por cero
entrega una señal escalonada y con ruido de período. Derivar una señal en
escalera produce picos bruscos de tensión cada vez que llega un cruce nuevo,
excitando innecesariamente el modo torsional del acople magnético. Para una
planta predominantemente de primer orden como un motor en velocidad, un PI es la
elección natural y robusta.

% ---------------------------------------------------------------------
\section{Diseño del PI a mano sobre \texorpdfstring{$G_1(s)$}{G1(s)}}
% ---------------------------------------------------------------------

Cerremos el lazo de $C_\omega(s) = \frac{k_p s + k_i}{s}$ con
$G_1(s) = \frac{K_{dc}}{\tau_{\text{mec}}s + 1}$ (despreciando primero el
retardo, que verificaremos a continuación):
\begin{equation}
  T(s) = \frac{C_\omega(s) G_1(s)}{1 + C_\omega(s) G_1(s)}
       = \frac{K_{dc}(k_p s + k_i)}{\tau_{\text{mec}} s^2 + (1 + K_{dc}k_p)s + K_{dc}k_i}.
  \label{eq:lazo_cerrado_pi}
\end{equation}
Dividiendo numerador y denominador por $\tau_{\text{mec}}$, el polinomio
característico del denominador tiene la forma canónica de segundo orden
$s^2 + 2\zeta\omega_n s + \omega_n^2$:
\[
  s^2 + \frac{1 + K_{dc}k_p}{\tau_{\text{mec}}}\,s + \frac{K_{dc}k_i}{\tau_{\text{mec}}}
  = s^2 + 2\zeta\omega_n\,s + \omega_n^2.
\]
Igualando coeficiente a coeficiente, despejamos las dos ganancias en dos líneas:
\begin{ecuacionclave}
\begin{equation}
  k_i = \frac{\tau_{\text{mec}}\,\omega_n^2}{K_{dc}},
  \qquad
  k_p = \frac{2\,\zeta\,\omega_n\,\tau_{\text{mec}} - 1}{K_{dc}}.
  \label{eq:sintonia_pi}
\end{equation}
\end{ecuacionclave}
\emph{Nota de unidades:} si en una implementación práctica el error se mide en
\si{rpm} en vez de \si{rad/s} y la salida en cuentas PWM sobre $u_{\max}$ en vez
de ciclo de trabajo $d\in[0,1]$, basta escalar $K_{dc}$ por el factor
constante $\frac{30}{\pi}\frac{1}{u_{\max}}$; la física del diseño es idéntica.

\paragraph{Aplicación al ejemplo numérico y chequeo de retardo.}
Con nuestra planta $G_1(s) = \frac{386}{0{,}118\,s + 1}$, elegimos una respuesta
suave y bien amortiguada ($\zeta = 0{,}85$) con ancho de banda moderado
$\omega_n = \SI{10}{rad/s}$ (más de 5 veces por debajo de la frecuencia
torsional $\omega_{n,\text{tors}} \approx \SI{54}{rad/s}$ para no sacudir la
cápsula!):
\[
  k_i = \frac{0{,}118\times 10^2}{386} \approx \SI{0,0306}{rad^{-1}},
  \qquad
  k_p = \frac{2(0{,}85)(10)(0{,}118) - 1}{386} \approx \SI{0,00261}{s/rad}.
\]
Ahora verifiquemos el efecto del retardo del sensor $\tau(n)$ (con $M=2$, $p=2$,
$\Delta t_c = \SI{2}{ms}$):
\begin{itemize}
  \item A $n_0 = \SI{1000}{rpm}$: $\tau(1000) \approx \SI{23,5}{ms}$, pérdida de
        fase a $\omega_c \approx \SI{10}{rad/s}$ de apenas
        $10\times 0{,}0235 = \SI{0,235}{rad} \approx \SI{13,5}{\degree}$
        (lazo bien estable).
  \item A $n = \SI{400}{rpm}$: $\tau(400) \approx \SI{57,3}{ms}$, la pérdida de
        fase sube a $\SI{33}{\degree}$. Si alguien hubiera diseñado un lazo
        agresivo con $\omega_c = \SI{25}{rad/s}$, a \SI{400}{rpm} el retardo por
        sí solo comería \SI{82}{\degree} de margen de fase y el motor empezaría
        a oscilar violentamente!
\end{itemize}

% ---------------------------------------------------------------------
\section{Discretización trapezoidal y saturación (\emph{anti-windup})}
% ---------------------------------------------------------------------

Un controlador digital no integra de forma continua, sino que actualiza su
salida cada $\Delta t_c$ segundos a partir de muestras $e_k$. La forma más limpia
y fiel de aproximar el área bajo la curva del error entre $t_{k-1}$ y $t_k$ es
por \textbf{trapecios} (método de Tustin):
\begin{equation}
  I_k = I_{k-1} + k_i\,\Delta t_c\,\frac{e_k + e_{k-1}}{2},
  \qquad
  d_k^{\text{lin}} = k_p\,e_k + I_k,
  \label{eq:tustin_cap8}
\end{equation}
cuya función de transferencia discreta en transformada $\mathcal{Z}$ es
$C(z) = k_p + \frac{k_i\Delta t_c}{2}\frac{z+1}{z-1}$.

\paragraph{El peligro del \emph{integrator windup}.}
Recordemos la saturación física $d_k \in [d_{\min}, 1]$ del
Capítulo~\ref{cap:motor}. Supongamos que pedimos un escalón grande o que el eje
se traba un instante: la salida $d_k$ pega contra el techo $d = 1$ (la fuente
está entregando todo su voltaje $V_{dc}$), pero la velocidad todavía está por
debajo de la consigna ($e_k > 0$). Si el integrador $I_k$ siguiera sumando error
mientras el actuador está saturado, acumularía un valor gigantesco
(\emph{windup} o «embalamiento integral»). Cuando finalmente el motor alcanza la
consigna, ese integrador inflado mantiene $d_k = 1$ durante demasiado tiempo y
produce un sobrepico enorme.

La solución física es inmediata (\emph{clamping}): \textbf{recortar el estado
integral} para que nunca empuje la salida más allá de lo que el inversor puede
entregar:
\begin{equation}
  I_k \leftarrow \sat\!\big(I_k,\; d_{\min} - k_p e_k,\; 1 - k_p e_k\big),
  \qquad
  d_k = \sat\!\big(k_p e_k + I_k,\; d_{\min},\; 1\big).
  \label{eq:antiwindup_cap8}
\end{equation}

% ---------------------------------------------------------------------
\section{El mapa del rango de validez: dónde acierta y dónde miente el modelo lineal}
% ---------------------------------------------------------------------

Si aplicamos el controlador PI diseñado sobre $G_1(s)$ a la \textbf{planta no
lineal completa}~\eqref{eq:nl_1}--\eqref{eq:nl_4} junto con el sensor BEMF y el
acople magnético, ambos coinciden muy bien dentro de una «isla segura» de
operación. Pero esa isla está rodeada por \textbf{cuatro paredes físicas} donde
el modelo lineal deja de representar la realidad. La
Figura~\ref{fig:mapa_validez} dibuja ese mapa en el plano velocidad--carga
$(n_0, T_{L,\text{tot}})$.

\begin{figure}[htbp]
\centering
\begin{tikzpicture}
\begin{axis}[
  width=12.0cm, height=6.8cm,
  axis lines=left,
  xlabel={Velocidad de operación $n_0$ [rpm]},
  ylabel={Torque de carga sobre la cápsula $T_{L,\text{tot}}$ [\si{mN.m}]},
  xmin=0, xmax=5200, ymin=0, ymax=9.0,
  xtick={350, 1000, 1600, 4775},
  xticklabels={$n_{\text{ciego}}$, $1000$, $1600$, $n_{\text{vacío}}$},
  ytick={2.89, 6.0},
  yticklabels={$2{,}89$, $T_{\max}=6$},
  tick label style={font=\footnotesize},
  label style={font=\small},
  clip=true
]
  % Zona 1: Ceguera del sensor (0 a 350 rpm)
  \fill[crimson!18] (axis cs:0,0) rectangle (axis cs:350,9.0);
  \draw[crimson, very thick] (axis cs:350,0) -- (axis cs:350,9.0);

  % Zona 2: Pérdida de margen de fase por retardo (350 a ~620 rpm)
  \fill[amberwarn!16] (axis cs:350,0) rectangle (axis cs:620,6.0);
  \draw[amberwarn, dashed, thick] (axis cs:620,0) -- (axis cs:620,6.0);

  % Zona 4: Techo del acople magnético (T > T_max = 6 mNm)
  \fill[crimson!12] (axis cs:350,6.0) rectangle (axis cs:5200,9.0);
  \draw[crimson, very thick] (axis cs:350,6.0) -- (axis cs:5200,6.0);

  % Zona 3: Techo de tensión d = 1: T_e(w) = 72 - 0.144*w [mNm],
  % restando T_f1 aprox 0.7 mNm -> llega a T=0 cerca de 4700 rpm y T=6 cerca de 4300 rpm
  % Pendiente en rpm: 0.144 * pi/30 = 0.01508 mNm/rpm
  % En n=4775 -> T_e=0; en n=4350 -> T_e = 6.4 mNm
  \fill[inkblue!14] (axis cs:4350,6.0) -- (axis cs:4740,0) -- (axis cs:5200,0) -- (axis cs:5200,6.0) -- cycle;
  \draw[inkblue, very thick] (axis cs:4180,9.0) -- (axis cs:4740,0);

  % Isla segura (válida para modelo lineal)
  \fill[tealaccent!10] (axis cs:620,0) -- (axis cs:4740,0) -- (axis cs:4350,6.0) -- (axis cs:620,6.0) -- cycle;

  % Curva de carga de agua del ejemplo: 0.3 + 9e-6*w + 1.5e-7*w^2 [Nm] -> en mNm:
  % w = pi/30 * n = 0.10472 * n
  % T = 0.3 + 0.0009425*n + 1.645e-6 * n^2
  \addplot[tealaccent, very thick, domain=0:2250, samples=80] {0.3 + 0.0009425*x + 0.000001645*x^2};

  % Punto de operación nominal (1000 rpm, 2.89 mNm)
  \fill[inkblue] (axis cs:1000,2.89) circle (3.2pt);
  \node[below right, inkblue, font=\scriptsize\bfseries] at (axis cs:1050,2.60) {Punto nominal};

  % Punto de cruce con T_max (~1600 rpm)
  \fill[crimson] (axis cs:1600,6.0) circle (3.0pt);
  \node[above right, crimson, font=\scriptsize\bfseries] at (axis cs:1680,6.12) {Desacople en agua};

  % Etiquetas de regiones
  \node[crimson, font=\scriptsize\bfseries, rotate=90] at (axis cs:175,4.2) {1. Zona ciega};
  \node[amberwarn, font=\scriptsize\bfseries, rotate=90] at (axis cs:485,3.0) {2. Retardo $1/n$};
  \node[crimson, font=\footnotesize\bfseries, fill=crimson!12, inner sep=2pt] at (axis cs:2350,7.75) {4. Desacople ($T_{L,\text{tot}} > T_{\max}$): la cápsula se frena};
  \node[tealaccent, font=\footnotesize\bfseries, align=center] at (axis cs:3080,2.85) {Ventana de validez\\del modelo lineal};
  \node[inkblue, font=\scriptsize\bfseries, rotate=-62] at (axis cs:4700,3.8) {3. Techo $d=1$};
  \node[tealaccent, font=\scriptsize, right, align=left] at (axis cs:1750,4.7) {Curva de carga\\(agua)};
\end{axis}
\end{tikzpicture}
\caption{Mapa de operación en el plano velocidad--carga ilustrado con los parámetros del ejemplo. El modelo linealizado y el PI diseñado a mano sólo son fieles dentro de la ventana verde delimitada por las cuatro paredes no lineales.}
\label{fig:mapa_validez}
\end{figure}

Repasemos brevemente qué ocurre al cruzar cada una de las cuatro fronteras de la
Figura~\ref{fig:mapa_validez}:
\begin{enumerate}
  \item \textbf{Pared izquierda --- Zona ciega ($n_1 < n_{\text{ciego}}$):} la
        BEMF desaparece bajo el ruido del comparador; el lazo cerrado pierde el
        \SI{100}{\percent} de su señal, dispara el timeout de \emph{stall} y
        obliga a reiniciar con la rampa a lazo abierto.
  \item \textbf{Franja de retardo ($n_{\text{ciego}} < n_1 < n_{\min,\text{estab}}$):}
        aunque hay cruces válidos, el período $T_{ZC}$ se alarga como $1/n$,
        roba margen de fase y vuelve oscilatorio un PI diseñado a alta velocidad.
  \item \textbf{Pared derecha --- Saturación de tensión ($d \to 1$ ó $d \to d_{\min}$):}
        el puente agota la tensión de bus $V_{dc}$; durante escalones grandes o
        cargas pesadas el lazo queda temporalmente abierto y sólo el
        \emph{anti-windup} evita sobrepicos violentos.
  \item \textbf{Techo superior --- Desacople magnético ($T_{L,\text{tot}} > T_{\max}$
        o salto dinámico $\delta > \delta_u$):} incluso con el motor bien lejos
        de saturar tensión, al subir la velocidad o agregar un líquido más
        viscoso la curva de carga cruza $T_{\max}$. El acople se ablanda
        ($K_{\text{eff}} \to 0$) y suelta la cápsula ($\omega_2 \to 0$), mientras
        el controlador sigue regulando $\omega_1 \approx n^\ast$ sin advertir el
        fallo.
\end{enumerate}

% ---------------------------------------------------------------------
\section{Síntesis final: toda la física del modelo en una sola hoja}
% ---------------------------------------------------------------------

Como cierre de este libro, la Tabla~\ref{tab:mapa_final} condensa de un vistazo
el hilo conductor completo: desde las cuatro leyes elementales hasta las
ecuaciones no lineales, su versión linealizada y las fronteras de validez. Si
podés recorrer cada fila de izquierda a derecha sobre una hoja en blanco con un
lápiz, \textbf{el modelo es completamente tuyo}.

\begin{table}[htbp]
\centering
\scriptsize
\setlength{\tabcolsep}{4pt}
\renewcommand{\arraystretch}{1.35}
\begin{tabular}{@{}
  >{\raggedright\arraybackslash}p{2.0cm}
  >{\raggedright\arraybackslash}p{2.1cm}
  >{\raggedright\arraybackslash}p{4.3cm}
  >{\raggedright\arraybackslash}p{3.4cm}
  >{\raggedright\arraybackslash}p{2.6cm}@{}}
\toprule
\textbf{Subsistema} & \textbf{Ley física} & \textbf{Ecuación no lineal} & \textbf{Aprox. lineal} & \textbf{Frontera de ruptura} \\
\midrule
Motor (six-step) & Kirchhoff + Faraday + Energía &
$2L_s\dot I = d V_{dc} - 2R I - 2k_e\omega_1$,\newline $T_e = 2k_e I$,\; $d\in[d_{\min},1]$ &
$2L_s\dot{\tilde I} = V_{dc}\tilde d - 2R\tilde I - 2k_e\tilde\omega_1$ (o $\tilde I$ cuasi-estática $\Rightarrow B_e=\frac{2k_e^2}{R}$) &
Saturación de tensión ($d=d_{\min}$ ó $d=1$) \\
\midrule
Rotor ($J_1$) & 2.ª y 3.ª ley de Newton &
$J_1\dot\omega_1 = T_e - T_c(\delta) - T_{f1}(\omega_1)$ &
$J_1\dot{\tilde\omega}_1 = 2k_e\tilde I - K_{\text{eff}}\tilde\delta - b_1\tilde\omega_1$ &
Desconexión cuando $\langle T_c\rangle \to 0$ \\
\midrule
Acople magnético & Cinemática / energía potencial &
$T_c(\delta) = T_{\max}\sin(N_a\delta)$,\newline $\dot\delta = \omega_1 - \omega_2$,\; $\delta_{\text{crit}}=\frac{\pi}{2N_a}$ &
$\tilde T_c = K_{\text{eff}}\,\tilde\delta$,\newline $K_{\text{eff}}=N_a\sqrt{T_{\max}^2-T_{L,\text{tot}}^2}$ &
$T_{L,\text{tot}} > T_{\max}$ o escalón que vence $\Delta U$ \\
\midrule
Agitador + fluido ($J_2$) & Newton + viscosidad + inercia &
$J_2\dot\omega_2 = T_c - T_L(\omega_2) - T_{f2}(\omega_2)$,\newline $T_L = b_{fl}\omega_2 + c_{fl}\omega_2|\omega_2|$ &
$J_2\dot{\tilde\omega}_2 = K_{\text{eff}}\tilde\delta - B_{2,\text{eq}}\tilde\omega_2$,\newline $B_{2,\text{eq}} = b_2 + b_{fl} + 2c_{fl}\omega_0$ &
Cambio de fluido / viscosidad que supera $T_{\max}$ \\
\midrule
Sensor BEMF + Arranque & Faraday (cruces por cero) &
$T_{ZC}=\frac{\pi}{p\omega_1}$, ciego si $n_1 < n_{\text{ciego}}$;\newline arranque síncrono $\propto\sin(\theta_s-p\theta_1)$ &
Retardo puro $e^{-s\tau(n_0)}$,\newline $\tau(n_0) \approx \frac{15(M+1)}{p\,n_0} + \frac{\Delta t_c}{2}$ &
Ceguera ($<n_{\text{ciego}}$) y pérdida de margen de fase ($\omega_c\tau\ge\Delta\phi$) \\
\bottomrule
\end{tabular}
\caption{Resumen integral del modelo físico, su linealización y las fronteras donde la aproximación lineal pierde validez.}
\label{tab:mapa_final}
\end{table}

\begin{lapizypapel}[Prueba final de dominio completo]
Imaginá que tenés cinco minutos de exposición y veinte minutos de preguntas con
lápiz y papel:
\begin{enumerate}[leftmargin=1.5em]
  \item ¿Podés explicar en 60 segundos, sin mirar apuntes, cuál es el problema
        físico completo y cuáles son las cuatro no linealidades más importantes?
  \item Dado $G_1(s) = \frac{K_{dc}}{\tau_{\text{mec}}s + 1}$, ¿podés derivar
        las fórmulas~\eqref{eq:sintonia_pi} de $k_p$ y $k_i$ y justificar por
        qué necesitás \emph{anti-windup}?
  \item Dibujá el mapa de la Figura~\ref{fig:mapa_validez} desde cero y explicá
        qué pasa físicamente cuando el sistema cruza cada una de sus cuatro
        paredes.
  \item Si alguien te pregunta \emph{«¿por qué no modelaste cogging, temperatura
        del bobinado ni tablas $N_p(\mathrm{Re})$?»}, ¿podés justificar en dos
        frases por qué dejarlos afuera vuelve al modelo más honesto y
        defendible?
\end{enumerate}
\end{lapizypapel}

\end{document}
